% Options for packages loaded elsewhere
\PassOptionsToPackage{unicode}{hyperref}
\PassOptionsToPackage{hyphens}{url}
\documentclass[
]{article}
\usepackage{xcolor}
\usepackage{amsmath,amssymb}
\setcounter{secnumdepth}{-\maxdimen} % remove section numbering
\usepackage{iftex}
\ifPDFTeX
  \usepackage[T1]{fontenc}
  \usepackage[utf8]{inputenc}
  \usepackage{textcomp} % provide euro and other symbols
\else % if luatex or xetex
  \usepackage{unicode-math} % this also loads fontspec
  \defaultfontfeatures{Scale=MatchLowercase}
  \defaultfontfeatures[\rmfamily]{Ligatures=TeX,Scale=1}
\fi
\usepackage{lmodern}
\ifPDFTeX\else
  % xetex/luatex font selection
  \setmainfont[]{Latin Modern Roman}
\fi
% Use upquote if available, for straight quotes in verbatim environments
\IfFileExists{upquote.sty}{\usepackage{upquote}}{}
\IfFileExists{microtype.sty}{% use microtype if available
  \usepackage[]{microtype}
  \UseMicrotypeSet[protrusion]{basicmath} % disable protrusion for tt fonts
}{}
\makeatletter
\@ifundefined{KOMAClassName}{% if non-KOMA class
  \IfFileExists{parskip.sty}{%
    \usepackage{parskip}
  }{% else
    \setlength{\parindent}{0pt}
    \setlength{\parskip}{6pt plus 2pt minus 1pt}}
}{% if KOMA class
  \KOMAoptions{parskip=half}}
\makeatother
\usepackage{longtable,booktabs,array}
\newcounter{none} % for unnumbered tables
\usepackage{calc} % for calculating minipage widths
% Correct order of tables after \paragraph or \subparagraph
\usepackage{etoolbox}
\makeatletter
\patchcmd\longtable{\par}{\if@noskipsec\mbox{}\fi\par}{}{}
\makeatother
% Allow footnotes in longtable head/foot
\IfFileExists{footnotehyper.sty}{\usepackage{footnotehyper}}{\usepackage{footnote}}
\makesavenoteenv{longtable}
\setlength{\emergencystretch}{3em} % prevent overfull lines
\providecommand{\tightlist}{%
  \setlength{\itemsep}{0pt}\setlength{\parskip}{0pt}}
\usepackage{setspace}
\onehalfspacing
\usepackage{bookmark}
\IfFileExists{xurl.sty}{\usepackage{xurl}}{} % add URL line breaks if available
\urlstyle{same}
\hypersetup{
  hidelinks,
  pdfcreator={LaTeX via pandoc}}

\author{}
\date{}

\begin{document}

\title{Prescribed Banzhaf Power by Exact Computation: Closed-Form Realizations and Certified Weighted Obstructions}
\author{Yongjin Chen\textsuperscript{a}, Shi Jin\textsuperscript{a}, Songze Zhu\textsuperscript{a,*}\\[4pt]
{\footnotesize \textsuperscript{a} School of Public Security Management, People's Public Security}\\[1pt]
{\footnotesize University of China, Beijing 100038, China}\\[2pt]
{\footnotesize \textsuperscript{*} Corresponding author. E-mail: 20196846@ppsuc.edu.cn}}
\date{}
\maketitle

\subsection{Abstract}\label{abstract}

Designing a voting rule that realizes a prescribed distribution of
voting power --- the inverse power index problem --- is a fundamental
computational task for parliaments, councils, and other weighted voting
bodies. This paper develops an exact, fully reproducible computational
methodology for the inverse Banzhaf problem and applies it to the
egalitarian benchmark target of Kurz (2015), in which a single player
holds half the voting power of each of the remaining players.

On the constructive side, a single parametric family of simple games
(the W-family) with an exact closed-form swing formula yields explicit
realizing games for every number of voters from six upward, resolving
Kurz's Conjecture 16 and correcting the sole erroneous published
example; a compact formulation with a cubic number of variables solves
every instance up to eighty voters in seconds to minutes. On the
obstruction side, a sharpened swap/block rigidity condition drives a
deterministic constraint-programming satisfiability (CP-SAT)
certification that no weighted majority game on six to ten voters
realizes the target, at a conclusive weight bound, and the certified
optimal weighted approximation is computed exactly for up to eight
voters.

As a computational case study, the methodology is applied to the
historical weighted Council of the European Union: its exact Banzhaf
power distributions are reproduced and shown to be far from the target,
with the certified optimal redesign for six voters improving the
distance about threefold. Every numerical claim is reproducible from the
companion code package.

\textbf{Keywords:} voting games; Banzhaf power index; inverse power
index problem; integer programming; constraint programming

\subsection{1. Introduction: the inverse Banzhaf
problem}\label{introduction-the-inverse-banzhaf-problem}

Designing a voting rule for a committee of \emph{dissimilar} members ---
a weighted parliament, the Council of the European Union, or a federal
upper chamber --- is a decision problem: the designer must decide how
voting \emph{power} is to be distributed among members. The Banzhaf
index (and the related Shapley--Shubik index) quantify power as the
normalized number of \emph{swings}, i.e., coalitions in which a member
is decisive (Penrose 1946; Felsenthal \& Machover 1998); for a weighted
majority game, computing these indices is NP-hard (Matsui \& Matsui
2001). The \textbf{inverse power index problem} is the formal dual of
this computation: given a desired power distribution, decide whether a
voting rule realizes it, construct one if so, and certify impossibility
otherwise (Nurmi 1982; Kurz \& Napel 2014; de Keijzer et al.~2014). From
an operations-research standpoint this is an \emph{inverse-optimization}
problem: the power index is the forward map from rules to power
distributions, and the designer must invert it subject to the structural
constraints of a chosen rule class --- general simple games, or the
weighted majority games used by most real voting bodies (Muroga 1971).

This paper contributes an exact, fully reproducible computational
methodology for both halves of this design decision problem, and applies
it to a benchmark target that has resisted a unified treatment since
2014: the vector \(\psi^n\) = (2, \(\dots,\) 2, 1) / \((2n-1),\) (1)
which assigns to \(n-1\) members twice the power of the remaining
member. Kurz (2015) conjectured the existence of a realizing simple game
for every n \(\ge\) 6:

\begin{quote}
\textbf{Conjecture 16 (Kurz 2015).} For every n \(\ge\) 6 there exists a
simple game v on n voters such that Bz(v) = \(\psi^n,\) where Bz is the
(normalized) Penrose--Banzhaf index.
\end{quote}

and the natural dual --- that no \emph{weighted majority} game can do so
(Kurz's Conjecture 15). Kurz and Napel (2014) verified Conjecture 16
computationally for 6 \(\le\) n \(\le\) 18, but gave no construction
valid for all n; we provide one (Corollaries 14--17, Section 11), and
settle the weighted obstruction computationally for n \(\le\) 10
(Section 9).

The paper is organized as an exact computational study, centered on two
engines. On the constructive side, Sections 3--7 give closed-form
families, exact swing formulas, and explicit games realizing \(\psi^n\)
for every n \(\ge\) 6; Section 8 summarizes the constructive half. On
the obstruction side, Section 9 derives a sharp structural rigidity
property \((M')\) and certifies by the constraint-programming
satisfiability (CP-SAT) solver that no weighted majority game realizes
\(\psi^n\) for n \(\le\) 10; Section 10 collects the two engines as a
reusable methodology, with complexity statements and instance tables;
Section 11 develops the unified W-family framework, the exact
closed-form constructions, and the open problems. Section 12 applies the
methodology as an exact computational case study to the historical
weighted Council of the European Union, whose actual Banzhaf power
distributions fall exactly in the certified range. Section 13 concludes.

\subsection{2. Preliminaries}\label{preliminaries}

Let N = \(\{1,\dots,n\}\) be a set of voters. A \textbf{simple game}
(Taylor \& Zwicker 1999) is a monotone Boolean function v :
\(2^N\ \to\ \{0,1\}\) with \(v(\varnothing)\) = 0 and v(N) = 1. A
coalition S \(\subseteq\) N is \textbf{winning} if v(S) = 1 and
\textbf{losing} otherwise. A winning coalition is \textbf{minimal
winning} if all its proper subsets are losing; a simple game is uniquely
determined by its family of minimal winning coalitions.

For i \(\in\) N, let \(\beta_i(v)\) =
\(\vert\{S\ \subseteq\ N\setminus\{i\}\) : v(S) = 0, v(S
\(\cup\ \{i\})\) = \(1\}\vert\) be the number of \textbf{swings} of
player i (the Banzhaf measure). The (normalized)
\textbf{Penrose--Banzhaf index} is \(\mathrm{Bz}_i(v)\) = \(\beta_i(v)\)
/ \(\sum_j\ \beta_j(v).\) Thus Bz(v) = \(\psi^n\) iff \(\beta(v)\) is
proportional to \((2,\dots,2,1),\) i.e., iff \(n-1\) players have an
equal number 2c of swings and one player has c swings, for some integer
c \(\ge\) 1.

\subsection{3. The complement-of-edges
family}\label{the-complement-of-edges-family}

For a graph G on vertex set N with edge set E, consider the simple game
\(v_G\) whose minimal winning coalitions are exactly the complements of
the edges of G: \(v_G(S)\) = 1 \(\iff\) S
\(\supseteq\ N\setminus\{a,b\}\) for some edge \(\{a,b\}\ \in\) E.
Equivalently, a coalition is winning iff it misses no edge of G. We
prove:

\textbf{Theorem 1.} Let G be a graph on n vertices with edge set E.
Suppose (i) G has no isolated vertices, and (ii) G is not a star (no
vertex is incident with every edge). Then the Banzhaf swing counts of
\(v_G\) are \(\beta_i(v_G)\) = \(\vert E\vert\) +
\((n-1)\ -\ 2\cdot\deg_G(i)\) (i \(\in\) N). (2)

\emph{Proof.} Player i has a swing on a losing coalition S with i
\(\notin\) S and S \(\cup\ \{i\}\) winning. The latter means S
\(\cup\ \{i\}\ \supseteq\ N\setminus\{a,b\}\) for some edge \(\{a,b\}.\)
- If \(\vert S\vert\) = \(n-3,\) then S \(\cup\ \{i\}\) =
\(N\setminus\{a,b\}\) is an \((n-2)-set,\) so S =
\(N(\{a,b\}\ \cup\ \{i\})\) with i \(\notin\ \{a,b\};\) such S is losing
since it has size \(n-3.\) Each edge not containing i contributes
exactly one such swing, giving \(\vert E\vert\ -\ \deg_G(i)\) swings. -
If \(\vert S\vert\) = \(n-2,\) then S \(\cup\ \{i\}\) =
\(N\setminus\{x\}\) for some x \(\ne\) i. This is winning iff
\(\deg_G(x)\ \ge\) 1, which holds by (i). The set S =
\(N\setminus\{x,i\}\) is losing iff \(\{x,i\}\ \notin\) E (otherwise S
contains the winning \((n-2)-set\ N\setminus\{x,i\}).\) Hence exactly
\((n-1)\ -\ \deg_G(i)\) swings of this type. - If \(\vert S\vert\) =
\(n-1,\) then S = \(N\setminus\{i\}\) and S \(\cup\ \{i\}\) = N is
winning. But \(N\setminus\{i\}\) is losing iff every edge contains i,
i.e., G is a star centered at i, excluded by (ii). Hence no swings of
this type. - No other cardinality yields a swing.

Summing gives \(\beta_i(v_G)\) = \(\vert E\vert\) +
\((n-1)\ -\ 2\cdot\deg_G(i).\ \blacksquare\)

\textbf{Corollary 2.} If G has all vertices but one of common degree d
and the exceptional vertex of degree d + c/2, and if \(\vert E\vert\) +
\((n-1)\ -\) 2d = 2c, then \(\mathrm{Bz}(v_G)\) = \(\psi^n.\) Since
\(2\vert E\vert\) = nd + c/2, the existence of such a graph is
equivalent to the solvability of 7c = \(d(2n-8)\) + \(4(n-1)\) (3) in
positive integers d, c with 0 \(\le\) d \(\le\ n-2\) and d + c/2
\(\le\ n-1,\) together with the graphicality of the degree sequence
\((d,\dots,d,\) d+c/2).

For 6 \(\le\) n \(\le\) 18, equation (3) admits such a solution exactly
for n \(\in\ \{6,\) 7, 9, 13, 14, 16, \(17\}.\) For n \(\equiv\) 4 (mod
7) there is no solution at all; for n \(\in\ \{8,\) 10, 12, \(15\}\) the
required exceptional degree exceeds \(n-1.\) Hence the
complement-of-edges family alone does \textbf{not} cover all n --- which
is consistent with the conjecture requiring a genuinely different
construction for the remaining values.

The condition (3) is solvable for \emph{infinitely many} n: for n
\(\le\) 30 the solvable values are 6, 7, 9, 13, 14, 16, 17, 20, 21, 23,
24, 27, 28, 30, \(\dots\) Indeed, the residue classes n \(\equiv\) 6, 0,
2, 3 (mod 7) are solvable for every admissible n with the
\emph{constant} common degree d = 2, 3, 1, 4 respectively, since then
\(\deg_G(n)\) = (4n+4)/7, (5n+7)/7, (3n+1)/7, (6n+10)/7 \(\le\ n-1.\)
Thus Conjecture 16 holds, by Theorem 1, for \textbf{four of the seven
residue classes} (i.e.~for infinitely many n); the values n \(\equiv\)
1, 4, 5 (mod 7) require a genuinely different construction (given in
Sections 4 and 9).

\subsection{4. The threshold-with-exceptions
family}\label{the-threshold-with-exceptions-family}

The complement-of-edges family does not cover every n.~For the remaining
small values a second, equally explicit family is useful. Let E be a set
of pairs and F a set of triples of {[}n{]}. Define the game v(E,F) by
v(S) = 1 \(\iff\ \vert S\vert\ \ge\ n-1,\) or \((\vert S\vert\) =
\(n-2\) and N\S \(\notin\) E), or \((\vert S\vert\) = \(n-3\) and
N\S \(\in\) F). (4) That is, the game is ``almost'' a majority-style
threshold at size \(n-2,\) with a set of \emph{holes} (N\S \(\in\) E)
removed at level \(n-2\) and a set of \emph{exceptions} (N\S \(\in\) F)
added at level \(n-3.\)

\textbf{Lemma (monotonicity).} v(E,F) is monotone (hence a simple game)
iff no edge of E is contained in a triple of F, i.e.~E \(\cap\) Pairs(F)
= \(\varnothing,\) where Pairs(F) is the set of all two-element subsets
of triples in F. Indeed, a winning \((n-3)-set\) \(N\setminus X\) (X
\(\in\) F) has a losing \((n-2)-superset\ N\setminus\{a,b\}\) exactly
when \(\{a,b\}\ \in\) E with \(\{a,b\}\ \subseteq\) X.

\textbf{Theorem 2.} For any E, F with E \(\cap\) Pairs(F) =
\(\varnothing,\) the swing counts of v(E,F) are \(\beta_i(v(E,F))\) =
\(C(n-1,2)\) + \(\vert F\vert\ -\ \vert E\vert\) +
\(2\cdot\deg_E(i)\ -\ 2\cdot r_i\) + \(e_i,\) (5) where \(\deg_E(i)\) is
the degree of i in the graph E, \(r_i\) = \(\vert\{X\ \in\) F : i
\(\in\ X\}\vert,\) and \(e_i\) = \(\vert\{X\ \in\) F : i \(\in\) X,
\(X\setminus\{i\}\ \in\ E\}\vert.\)

\emph{Proof sketch.} Player i swings on a losing S with \(S\cup\{i\}\)
winning. Enumerate by \(\vert S\vert:\) - \(\vert S\vert\) = \(n-2\) (S
losing means N\S \(\in\) E): the unique losing \((n-2)-set\) with i
\(\notin\) S and N\S \(\in\) E requires N\S \(\ni\) i, giving
\(\deg_E(i)\) swings; - \(\vert S\vert\) = \(n-3:\) S =
\(N\setminus\{i,a,b\}\) with \(\{i,a,b\}\ \notin\) F and
\(\{a,b\}\ \notin\) E, giving \(C(n-1,2)\ -\ \vert E\vert\) +
\(\deg_E(i)\ -\ r_i\) + \(e_i\) swings; - \(\vert S\vert\) = \(n-4:\) S
= \(N(X\cup\{i\})\) for X \(\in\) F with i \(\notin\) X, giving
\(\vert F\vert\ -\ r_i\) swings. Summing yields (5). \(\blacksquare\)
(The identity is verified exhaustively against direct computation on
random instances for n = 8, 10, 12.)

\textbf{Corollary 3.} If \(\deg_E(i)\ -\ r_i\) is constant, say D, for
all i \(\ne\) n and \(\deg_E(n)\ -\ r_n\) = D \(-\) c/2, then Bz(v(E,F))
= \(\psi^n\) with scale c = \((n\cdot C(n-1,2)\) +
\((n-6)\vert F\vert\ -\ (n-4)\vert E\vert)/(2n-1)\) (when this is a
positive integer). In particular, such games realize \(\psi^n\) for
every n \(\in\ \{8,\) 10, 11, 12, \(15\}\) not covered by Corollary 2
(found by solving (5) via mixed-integer programming; see Section 7),
while for n = 18 a solution was obtained by Kurz and Napel.

\subsection{\texorpdfstring{5. Explicit games realizing
\(\psi^n\)}{5. Explicit games realizing \textbackslash psi\^{}n}}\label{explicit-games-realizing-ux3c8ux207f}

For each n \(\in\ \{6,7,9,13,14,16,17\}\) we give the graph \(G_n\)
(edges written as pairs) such that \(\mathrm{Bz}(v_{G_n})\) =
\(\psi^n;\) the minimal winning coalitions of \(v_{G_n}\) are the
complements of these edges. The resulting swing vector is \(\beta\) =
\(c\cdot(2,\dots,2,1),\) displayed as a multiple of c.

{\def\LTcaptype{none} % do not increment counter
\begin{longtable}[]{@{}
  >{\raggedright\arraybackslash}p{(\linewidth - 6\tabcolsep) * \real{0.1053}}
  >{\raggedright\arraybackslash}p{(\linewidth - 6\tabcolsep) * \real{0.5789}}
  >{\raggedright\arraybackslash}p{(\linewidth - 6\tabcolsep) * \real{0.1053}}
  >{\raggedright\arraybackslash}p{(\linewidth - 6\tabcolsep) * \real{0.2105}}@{}}
\toprule\noalign{}
\begin{minipage}[b]{\linewidth}\raggedright
n
\end{minipage} & \begin{minipage}[b]{\linewidth}\raggedright
graph \(G_n\)
\end{minipage} & \begin{minipage}[b]{\linewidth}\raggedright
c
\end{minipage} & \begin{minipage}[b]{\linewidth}\raggedright
\(\beta(v_{G_n})\)
\end{minipage} \\
\midrule\noalign{}
\endhead
\bottomrule\noalign{}
\endlastfoot
6 & \(C_4(1-2-6-3-1)\ \cup\ K_3(4,5,6)\) & 4 & (8,8,8,8,8,4) \\
7 & \(C_6\) on \(\{1,\dots,6\}\ \cup\) star centered at 7 & 6 &
\((12,\dots,12,6)\) \\
9 & star at 9 on \(\{1,2,3,4\}\ \cup\) matching
\(\{5,6\},\allowbreak\{7,8\}\) & 6 & \((12,\dots,12,6)\) \\
13 & star at 13 on \(\{1,\dots,8\}\ \cup\) matching on
\(\{1,\dots,8\}\ \cup\ C_4\) on \(\{9,\dots,12\}\) & 12 &
\((24,\dots,24,12)\) \\
14 & \(C_{13}\) on \(\{1,\dots,13\}\) + chord \(\{6,10\}\ \cup\) star at
14 except \(\{6,10\}\) & 16 & \((32,\dots,32,16)\) \\
16 & star at 16 on \(\{1,\dots,7\}\ \cup\) matching on
\(\{8,\dots,15\}\) & 12 & \((24,\dots,24,12)\) \\
17 & star at 17 (to all) \(\cup\ C_{16}\) on \(\{1,\dots,16\}\ \cup\)
matching \(\{i,i+8\}\ (1\le i\le8)\) & 24 & \((48,\dots,48,24)\) \\
\end{longtable}
}

Each entry is verified directly by computing the swing counts. For
instance, the n = 6 entry is the game whose minimal winning coalitions
are
\(\{1,2,3,4\},\allowbreak\ \{1,2,3,5\},\allowbreak\ \{1,2,3,6\},\allowbreak\ \{1,2,4,5\},\allowbreak\ \{1,3,4,5\},\allowbreak\ \{2,4,5,6\},\allowbreak\ \{3,4,5,6\}\)
(the complements of the seven edges of \(C_4\ \cup\ K_3),\) whose swing
counts are (8,8,8,8,8,4).

\subsection{6. A correction to the published
example}\label{a-correction-to-the-published-example}

Kurz and Napel (2014, footnote 19) state that the game on six voters
whose minimal winning coalitions are
\(\{2,4,5,6\},\allowbreak\ \{2,3,4,5\},\allowbreak\ \{1,3,5,6\},\allowbreak\ \{1,3,4,5\},\allowbreak\ \{1,2,4,6\},\allowbreak\ \{1,2,3,5\}\)
attains the PBI vector (8,8,8,8,8,4)/44 = \(\psi^6.\) This is
\textbf{incorrect}: a direct computation (or a routine check by hand
for, e.g., player 6, who has five swings) shows the swing counts are
\(\beta\) = (7,7,7,7,9,5), \(\sum\beta\) = 42, so that Bz(v) =
(7,7,7,7,9,5)/42 \(\ne\ \psi^6.\) The claim that an exact solution
exists for n = 6 is nevertheless true: the game in Section 5 (seven
four-sets) attains \(\psi^6.\) (An exhaustive enumeration of all
7,828,354 simple games on six voters yields 360 games, in two symmetry
classes, whose Banzhaf index equals \(\psi^6.)\)

\subsection{\texorpdfstring{7. Computational verification for 6 \(\le\)
n \(\le\)
17}{7. Computational verification for 6 \textbackslash le n \textbackslash le 17}}\label{computational-verification-for-6-n-17}

For the values n \(\in\ \{8,\) 10, 11, 12, \(15\}\) not covered by
Corollary 2 we verified existence by formulating the problem as a
mixed-integer program: variables x(S) \(\in\ \{0,1\}\) for the winning
status of each coalition S, the monotonicity constraints x(S) \(\le\)
x(T) for S \(\subseteq\) T, the normalization \(x(\varnothing)\) = 0,
x(N) = 1, and the swing equations
\(\sum_{S\ \not\ni\ i}\ (x(S\cup\{i\})\ -\) x(S)) = 2c (i \(\ne\) n),
\(\sum_{S\ \not\ni\ n}\ (x(S\cup\{n\})\ -\) x(S)) = c, with c \(\ge\) 1
integer. Using HiGHS, all these instances are feasible. Table 2
summarizes the scale c found (n = 6, 7 by exhaustive enumeration; n =
13, 14, 16, 17 by the closed-form construction of Section 4; the five
remaining values by mixed-integer linear programming (MILP) --- since
reproduced at the same scales by the closed-form \(\deg_G\) + \(r_i\) =
D constructions of Corollary 16).

{\def\LTcaptype{none} % do not increment counter
\begin{longtable}[]{@{}lllllllllllll@{}}
\toprule\noalign{}
n & 6 & 7 & 8 & 9 & 10 & 11 & 12 & 13 & 14 & 15 & 16 & 17 \\
\midrule\noalign{}
\endhead
\bottomrule\noalign{}
\endlastfoot
c & 4 & 6 & 8 & 6 & 10 & 14 & 16 & 12 & 16 & 20 & 12 & 24 \\
\end{longtable}
}

\emph{Note.} The scale c = 20 for n = 15 was found by the CP-SAT model
--- an improvement over the previously reported c = 22; minimality at n
= 15 is open (a downward sweep at c = 18 did not terminate within a 600
s solver budget). Minimality is certified for n = 8 (c = 8, with c
\(\in\ \{2,\) 4, \(6\}\) proven infeasible), n = 9 (c = 6, c
\(\in\ \{2,\ 4\}\) infeasible) and n = 10 (c = 10, c \(\in\ \{2,\) 4, 6,
\(8\}\) infeasible); see Open problem 3 in Section 11.

For n = 6 the existence is established by exhaustive enumeration over
all 7,828,354 simple games (360 solutions in two symmetry classes). For
n = 7, exhaustive enumeration of all games is infeasible; we instead
enumerated the 5-uniform family (winning coalitions of size 5)
exhaustively and found 490 solutions, all of the complement-of-edges
type \(C_6\ \cup\ star_7.\) Thus Conjecture 16 holds for every 6 \(\le\)
n \(\le\) 17 by the constructions of Sections 3--5; the value n = 18 and
all larger values are covered by the closed-form constructions of
Section 11 (Corollaries 14--17).

\subsection{8. Summary of the constructive
half}\label{summary-of-the-constructive-half}

For the inverse voting-design problem, we have delivered the
\emph{construction half} in full: an explicit, closed-form simple game
whose Penrose--Banzhaf index equals \(\psi^n\) =
\((2,\dots,2,1)/(2n-1)\) for every n \(\ge\) 6 --- constructively
resolving Kurz's Conjecture 16, whose benchmark status motivated the
design study. The main ingredients are:

\begin{enumerate}
\def\labelenumi{\arabic{enumi}.}
\tightlist
\item
  \textbf{Two closed-form families} with exact swing formulas (Theorems
  1 and 2) --- the \emph{complement-of-edges} family (covering n
  \(\equiv\) 0, 2, 3, 6 (mod 7)) and the \emph{triangle-with-special-
  triples} family (covering n \(\equiv\) 1, 4, 5 (mod 7)) --- unified in
  a single parametric framework, the W-family (Proposition 11), whose
  swing formula (U) is exact;
\item
  \textbf{A correction of the only published example} (Kurz and Napel
  2014, footnote 19);
\item
  \textbf{A complete closed-form construction for every n \(\ge\) 6}
  (Corollaries 14, 16, 17), where the five smallest values 8, 10, 11,
  12, 15 are handled by the \(\deg_G\) + \(r_i\) = D structure.
\end{enumerate}

The arithmetic characterization (3) shows that no single
complement-of-edges construction can cover all n, and the residue
classes differ genuinely; nevertheless, a single parametric framework
--- the W-family with base degree and a suitable number of (special)
triangles --- realizes \(\psi^n\) for every n.

\subsection{\texorpdfstring{9. The weighted case: \(\psi^n\) is not
realizable by any weighted majority
game}{9. The weighted case: \textbackslash psi\^{}n is not realizable by any weighted majority game}}\label{the-weighted-case-ux3c8ux207f-is-not-realizable-by-any-weighted-majority-game}

We now turn to the \emph{obstruction half} of the design decision
problem: is the target \(\psi^n\) realizable within the weighted
majority class --- the rule class used by most real voting bodies
(parliaments, the EU Council, the IMF, corporate voting)? This is the
question a designer must answer before deciding whether to stay in the
weighted class. It is precisely Kurz's \textbf{Conjecture 15}, the
natural ``dual'' of Conjecture 16. As stated by Kurz (2015, Conjecture
15) it is quantitative: there exists a universal c \textgreater{} 0 with
\(\|\mathrm{Bz}(v)\ -\ \psi^n\|_1\ \ge\) c/n for every n \(\ge\) 2 and
every weighted majority game v. In particular no weighted majority game
reproduces \(\psi^n\) exactly (distance zero is excluded). We certify
the exact case for n \(\le\) 10, and the quantitative form for n \(\le\)
8. Recall that a \textbf{weighted majority game} is a simple game {[}q;
\(w_1,\dots,w_n]\) with integer weights 1
\(\le\ w_1\ \le\ \dots\ \le\ w_n\) and quota q \(\ge\) 1, where S is
winning iff W(S) = \(\sum_{i\in S}\ w_i\ \ge\) q. Let \(\eta_i\) =
\(\beta_i([q;w])\) be the swing count. The following identity is a
classical structural fact for weighted majority games --- a swap/block
argument of the type long used in the Banzhaf literature (cf. Felsenthal
\& Machover 1998). We do not claim it as new; what is new is that it
rigidly dissects Kurz's specific target \(\psi^n,\) pinning the
obstruction on a subset-sum problem.

\textbf{Proposition 4 (structural decomposition; swap/block argument).}
Let v = {[}q; w{]} be a weighted majority game with \(w_i\ \ge\ w_j.\)
Then \(\eta_i\ -\ \eta_j\) = 2 \(\cdot\ \vert\{\) T
\(\subseteq\ N\setminus\{i,j\}\) : q \(-\ w_i\ \le\) W(T) \textless{} q
\(-\ w_j\ \}\vert.\) (K)

\emph{Proof.} Enumerate the swings of i and j by the subset T
\(\subseteq\ N\setminus\{i,j\}\) that excludes both, and by whether the
swing coalition S does or does not contain the other player: - i swings
on S = T iff W(T) \(\in\ [q-w_i,\) q); - i swings on S = T
\(\cup\ \{j\}\) iff W(T) \(\in\ [q-w_i-w_j,\ q-w_j);\) - j swings on S =
T iff W(T) \(\in\ [q-w_j,\) q); - j swings on S = T \(\cup\ \{i\}\) iff
W(T) \(\in\ [q-w_i-w_j,\ q-w_i).\)

Because \(w_i\ \ge\ w_j,\) we have \(q-w_i\ \le\ q-w_j\) and
\(q-w_i-w_j\ \le\ q-w_i,\) so \([q-w_j,\) q) \(\subseteq\ [q-w_i,\) q)
and \([q-w_i-w_j,\ q-w_i)\ \subseteq\ [q-w_i-w_j,\ q-w_j),\) and in both
inclusions the \emph{removed} interval is the same one,
\([q-w_i,\ q-w_j).\) Summing the four counts over all T and subtracting
\(\eta_j\) from \(\eta_i\) cancels the two common pieces and leaves
exactly twice the count of T with W(T)
\(\in\ [q-w_i,\ q-w_j).\ \blacksquare\)

The identity was verified on 12,000 random weighted games (n =
\(3\dots12,\) weights \(\le\) 30) with 0 failures
(\texttt{probe6\_\allowbreak identity.\allowbreak py}).

\textbf{Corollary 5 (parity; resolves Observation 7 in the Appendix).}
If Bz(v) = \(\psi^n,\) then the scale factor c is even. Two independent
proofs: - \emph{(any simple game, Dubey--Shapley 1979)}: in every simple
game all swing counts have a common parity. Here \(\beta\) =
\((2c,\dots,2c,c);\) since 2c is even, the common parity is even, so c
is even. - \emph{(weighted, from (K))}: if \(w_s\) is the special player
\((\eta_s\) = c), then for any i \(\ne\) s, c = \(\eta_i\ -\ \eta_s\) =
\(2\cdot\vert\{T\) : \(q-w_i\ \le\) W(T) \textless{} \(q-w_s\}\vert,\)
so c is even.

\textbf{Corollary 6 (the special player has strictly minimal weight).}
Suppose \(\eta_s\) is the unique minimum swing count and \(\eta_i\)
\textgreater{} \(\eta_s\) for every i \(\ne\) s. Then \(w_s\)
\textless{} \(w_i\) for all i \(\ne\) s. Indeed \(\eta_i\ -\ \eta_s\)
\textgreater{} 0 forces the interval \([q-w_i,\ q-w_s)\) to contain some
subset sum, so it is nonempty and \(w_i\) \textgreater{} \(w_s.\) Thus,
in any weighted game whose Banzhaf index is proportional to \(\psi^n,\)
the singled-out voter must be the \emph{unique lightest} voter --- an
a-priori monotonicity constraint on the weights.

\textbf{Corollary 7 (interval decomposition and the mirror
obstruction).} With s the special player and \(\eta\) =
\((2c,\dots,2c,c),\) identity (K) applied to the pair (i, s) yields, for
every i \(\ne\) s, \(\vert\{\) T \(\subseteq\ N\setminus\{i,s\}\) : q
\(-\ w_i\ \le\) W(T) \textless{} q \(-\ w_s\ \}\vert\) = c/2, (L) so the
excess c/2 of each non-special voter over the special voter is realized
by subset sums of \(N\setminus\{i,s\}\) inside the weight gap
\([q-w_i,\ q-w_s).\) Now fix two non-special voters i, j with
\(w_i\ \ge\ w_j\) (WLOG) and split the window count in (L) at (i, s)
according to whether j lies in T, and likewise the count at (j, s)
according to whether i lies in T:
\(\vert\{T\ \subseteq\ N\setminus\{i,s\}\) : \(q-w_i\ \le\) W(T)
\textless{} \(q-w_s\}\vert\) =
\(\vert\{T\ \subseteq\ N\setminus\{i,j,s\}\) : \(q-w_i\ \le\) W(T)
\textless{} \(q-w_s\}\vert\) +
\(\vert\{T\ \subseteq\ N\setminus\{i,j,s\}\) : \(q-w_i-w_j\ \le\) W(T)
\textless{} \(q-w_s-w_j\}\vert,\)
\(\vert\{T\ \subseteq\ N\setminus\{j,s\}\) : \(q-w_j\ \le\) W(T)
\textless{} \(q-w_s\}\vert\) =
\(\vert\{T\ \subseteq\ N\setminus\{i,j,s\}\) : \(q-w_j\ \le\) W(T)
\textless{} \(q-w_s\}\vert\) +
\(\vert\{T\ \subseteq\ N\setminus\{i,j,s\}\) : \(q-w_i-w_j\ \le\) W(T)
\textless{} \(q-w_s-w_i\}\vert,\) both equal to c/2. Since
\(w_i\ \ge\ w_j\) \textgreater{} \(w_s\) (Corollary 6), the relevant
weight windows nest as disjoint unions \([q-w_i,\ q-w_s)\) =
\([q-w_i,\ q-w_j)\ \sqcup\ [q-w_j,\ q-w_s),\)
\([q-w_i-w_j,\ q-w_s-w_j)\) =
\([q-w_i-w_j,\ q-w_s-w_i)\ \sqcup\ [q-w_s-w_i,\ q-w_s-w_j),\) so
subtracting the two equalities leaves the \textbf{mirror identity}
\(\vert\{\) T \(\subseteq\ N\setminus\{i,j,s\}\) : \(q-w_i\ \le\) W(T)
\textless{} \(q-w_j\ \}\vert\) + \(\vert\{\) T
\(\subseteq\ N\setminus\{i,j,s\}\) : \(q-w_s-w_i\ \le\) W(T) \textless{}
\(q-w_s-w_j\ \}\vert\) = 0. (M) Both summands are nonnegative, hence
they vanish individually: \(\vert\{\) T
\(\subseteq\ N\setminus\{i,j,s\}\) : \(q-w_i\ \le\) W(T) \textless{}
\(q-w_j\ \}\vert\) = 0 and \(\vert\{\) T
\(\subseteq\ N\setminus\{i,j,s\}\) : \(q-w_s-w_i\ \le\) W(T) \textless{}
\(q-w_s-w_j\ \}\vert\) = 0. \((M')\) In words: for every pair of
non-special voters i, j with \(w_i\ \ge\ w_j,\) the weight-gap window
\([q-w_i,\ q-w_j)\) --- and its translate by \(w_s\) --- is
\emph{sum-free} for the profile of the remaining \(n-3\) voters.
Together with (L) this is a sharp rigidity statement: the c/2 sums
realizing voter i's excess over s must come exclusively from the part
\([q-w_i-w_j,\ q-w_s-w_j)\) of the gap window, never from the sub-window
\([q-w_i,\ q-w_j)\) itself; and the reflected sub-window
\([q-w_s-w_i,\ q-w_s-w_j)\) is likewise empty. (Verified numerically:
across \(2.28\times10^5\) random weighted games with \(\eta_i\) =
\(\eta_j,\) both windows in \((M')\) were empty in every case.)

\textbf{The Banzhaf--weighted obstruction property \((M').\)} We record
the rigidity statement just derived as the paper's central structural
obstruction. Under the exact hypothesis \(\eta\) = \((2c,\dots,2c,c),\)
\emph{any} realizing game must satisfy (L) and \((M')\) --- they are
necessary conditions. They are also so tight that the full constraint
system --- (L), \((M'),\) together with the exact swing profile --- is
\emph{provably infeasible}, as certified by CP-SAT at the classical
weight bound \(w_i\ \le\ 2^{n-1}\) for n = 6, \(\dots,\) 10 (Section 9,
below). In this sense \((M')\) is the structural certificate of the
obstruction: the subset-sum windows it forces to be empty are precisely
the degrees of freedom whose infeasibility the deterministic solver
establishes. The property is not merely cosmetic: the mirror constraints
of \((M')\) are the pruning constraints of the enhanced feasibility
model (\texttt{conj15\_\allowbreak feas2.\allowbreak py}), which decides
n = 6 in 2.5 s (vs 2.9 s) and n = 7 in 19 s --- about six times faster
than the unpruned model at n = 7 (117 s) --- and are thus what makes the
n \(\ge\) 9 frontier approachable at all. We are not aware of this
reflected sum-freeness rigidity in the threshold-logic or voting-power
literature; the closest structural studies --- Chow's (1961) theorem on
Chow parameters (a recovery problem) and the inverse-power-index
certification results of Alon \& Edelman (2010) and Kurz (2016) ---
address reconstruction and worst-case approximation, not the
exact-hypothesis rigidity used here.

\textbf{Two-tier generalization \(\psi^n(a,b)\) and a family-level
no-go.} For a = b the target \(\psi^n(a,b)\) =
\((a,\dots,a,b)/((n-1)a+b)\) equals \(\psi^n\) with c scaled. For a
\textgreater{} b the W-family triangle construction of Proposition 11
realizes \(\psi^n(a,b)\) with special degree s = {[}
\((a-b)(5n-5-p)\ -\) 4aq + 12b {]} / (3a+b), (N) where p pure and q
special triangles are used. Systematic search (probe4/probe5) gives: -
\textbf{(3,1) is uniform}: for all n \(\in\) {[}6, 80{]}, s = \(n-1\)
with p = 0, q = 1 realizes \(\psi^n(3,1)\) --- the construction G =
star(n) \(\cup\ C_{n-1}\ \cup\ \{1,2\},\allowbreak\) F = \(\{triangle\)
on the special \(player\}\) works for every n, no case distinction. -
\textbf{(k,1), k \(\ge\) 4 is obstructed within this family}: for k
\(\ge\) 4 the required degree s \(\approx\ 5(k-1)n/(3k+1)\) exceeds
\(n-1\) for large n, so no triangle construction of this shape exists.
Coverage rates over n \(\in\) {[}6,80{]}: k=2: 73/75 (gaps 8, 11); k=3:
75/75; k=4: 40/75; k=5: 28/75; k=6: 9/75; k=7: 8/75. Combined with
Corollaries 6--7, this is strong evidence that \(\psi^n(k,1)\) for k
\(\ge\) 4 is not weighted-realizable at all; a complete proof is the
subject of ongoing work (Corollary 6 pins \(w_s\) minimal; (L) and (M)
then over-constrain the subset-sum profile).

\textbf{Computational obstruction (exact).} For \(\psi^n\) itself, exact
feasibility of a weighted realization with swing vector
\((2c,\dots,2c,c)\) was settled by CP-SAT
(\texttt{conj15\_\allowbreak feas.\allowbreak py}). Infeasibility is
certified at the classical weight bound W = \(2^{n-1}:\) every weighted
majority game on n voters has an integer representation with all weights
in {[}1, \(2^{n-1}]\) (a standard bound, implied for n \(\ge\) 3 already
by Muroga's sharper minimal-weight bound \(3^{n/3}),\) so infeasibility
at W = \(2^{n-1}\) is conclusive:

{\def\LTcaptype{none} % do not increment counter
\begin{longtable}[]{@{}lllll@{}}
\toprule\noalign{}
n & W = \(2^{n-1}\) & model & verdict & runtime \\
\midrule\noalign{}
\endhead
\bottomrule\noalign{}
\endlastfoot
6 & 32 & standard & INFEASIBLE & 2.9 s \\
7 & 64 & standard & INFEASIBLE & 117 s \\
8 & 128 & +symmetry breaking & INFEASIBLE & 316 s \\
9 & 256 & enhanced (\texttt{feas2}) & INFEASIBLE & 1562.9 s \\
10 & 512 & enhanced (\texttt{feas2}) & INFEASIBLE & 9527.4 s \\
\end{longtable}
}

The pipeline is three-tiered, and we are explicit about which tier backs
each claim. \textbf{Tier 1 (exhaustive)}: n = 6 is settled by
enumerating all 1,111 weighted games, guaranteed. \textbf{Tier 2
(deterministic)}: n = 7, 8, 9 and 10 are settled by CP-SAT at the
conclusive weight bound W = \(2^{n-1}\) (standard model for n = 7, 8;
the enhanced model below for n = 9, 10), guaranteed. \textbf{Tier 3
(budgeted)}: n \(\ge\) 11 runs the same solver under a time budget, so a
negative outcome there is heuristic evidence for a trend, not a proof,
and is labeled accordingly.

n = 6 agrees with the exhaustive enumeration of all 1,111 weighted games
on 6 voters (\texttt{weighted\_\allowbreak enum.\allowbreak py}), whose
best \(L_1-distance\) to \(\psi^6\) is 5/44 \(\approx\) 0.1136; hence no
weighted game attains distance 0. The n = 8 run required the
symmetry-breaking constraint that non-special voters can be relabeled in
nondecreasing weight order (all have equal target swing 2c); with it the
infeasibility certificate is found in 316 s, whereas the unconstrained
model is inconclusive after 20 min. n = 9 is substantially harder: the
standard model exhausts a full 2 h budget without a certificate
(UNKNOWN), a dramatic jump consistent with the exponential growth of the
subset-sum constraint structure. The enhanced model
(\texttt{conj15\_\allowbreak feas2.\allowbreak py}), which additionally
enforces the sharp mirror constraints \((M')\) --- valid since they are
consequences of the exact hypothesis --- then decides n = 9 in 1562.9 s
(3584 mirror constraints): \textbf{INFEASIBLE}. At n = 10 the same
enhanced model decides \textbf{INFEASIBLE} in 9527.4 s (9216 mirror
constraints), while the standard model again exhausts a 2 h budget
without a certificate (UNKNOWN). This closes the exact obstruction for
all n \(\le\) 10, the strongest conclusive statement currently
available; the enhanced model already gives a \(6\times\) speedup on the
smaller scales (n = 6: 2.5 s, n = 7: 19 s), and its growing
mirror-constraint set is what makes the n \(\ge\) 9 frontier
approachable at all. This is exactly Kurz's Conjecture 15, which asserts
the quantitative statement that for some universal c \textgreater{} 0,
\(\|\mathrm{Bz}(v)\ -\ \psi^n\|_1\ \ge\) c/n for \emph{every} weighted
majority game v and every n \(\ge\) 2. The exact obstruction for n = 6,
\(\dots,\) 10 is the first step; the constant c is extracted from the
optimal \(L_1\) errors computed just below.

\textbf{Optimal \(L_1-distance\) (exact computation,
\texttt{conj15\_\allowbreak l1.\allowbreak py}).} For each n we compute
the exact minimal value of \(\|\mathrm{Bz}(v)\ -\ \psi^n\|_1\) over all
weighted majority games, by fixing the total number of swings T =
\(\sum\eta_i\) and minimizing the integer objective
\(\sum_{i\ne s}\ \vert(2n-1)\eta_i\ -\ 2T\vert\) +
\(\vert(2n-1)\eta_s\ -\ T\vert\) in CP-SAT, enumerating T over its full
range and using the same weight bound W = \(2^{n-1};\) the special
player is fixed at position n without loss (relabeling). The method
reproduces the exhaustive n = 6 optimum exactly
(\texttt{conj15\_\allowbreak l1.\allowbreak py} n=6 \(T\le45):\)
\textbf{\(L_1\) = 5/44 \(\approx\) 0.113636 at T = 32, with weights w =
(3,1,3,3,3,1), quota q = 11}, matching the brute-force enumeration of
all 1,111 weighted games --- a strong check of both methods.

{\def\LTcaptype{none} % do not increment counter
\begin{longtable}[]{@{}lllll@{}}
\toprule\noalign{}
n & optimal \(L_1\) & T & \(n\cdot L_1\) & 1/n \\
\midrule\noalign{}
\endhead
\bottomrule\noalign{}
\endlastfoot
6 & 5/44 \(\approx\) 0.113636 & 32 & \(\approx\) 0.682 & 0.167 \\
7 & \(\approx\) 0.085470 & 45 & \(\approx\) 0.598 & 0.143 \\
8 & \(\approx\) 0.066667 & 60 & \(\approx\) 0.533 & 0.125 \\
\end{longtable}
}

(The n = 8 value is certified as the global optimum: T = 60 is OPTIMAL
by CP-SAT, and the total swing count was swept over its full range to
the theoretical maximum \(8\cdot2^7\) = 1024 ---
\texttt{conj15\_\allowbreak l1\_\allowbreak range.\allowbreak py} ---
whose best value outside T = 60 is \(L_1\) = 0.116667 at T = 280, well
above the incumbent.) For n = 7 the optimal profile is (5,7,7,7,7,7,5)
with weights (2,64,64,64,64,64,1) and quota 65; the value is certified
OPTIMAL by CP-SAT (re-verified with a 600 s budget), and the full sweep
over T \(\le\) 448 --- the theoretical maximum \(7\cdot2^6\) of the
total swing count --- confirms that no other total swing count improves
on it: the best value in the complementary range T \(\ge\) 121 is
\(L_1\) = 0.119535 at T = 139
(\texttt{conj15\_\allowbreak l1\_\allowbreak range.\allowbreak py}),
well above the incumbent. For n = 8 the optimum is attained at profile
(6,8,8,8,8,8,8,6) with weights (2,3,3,3,3,3,3,1) and quota 4. The column
\(n\cdot L_1\) is the empirical anchor for the universal constant in
Conjecture 15: since \(\|\cdot\|_1\ \ge\) c/n, any admissible c is at
most \(\inf_n\ n\cdot L_1(n);\) the values above give c \(\le\)
min(30/44, \(7\cdot0.08547,\ 8\cdot0.066667)\ \approx\) 0.533. We stress
the distinction between proof and certification, which is central to how
the computational results below are to be read: the value 5/44 for n = 6
is an exhaustive optimum, and 0.085470 for n = 7 is a \emph{certified}
optimum (over the full range T \(\le\) 448 of the total swing count);
the value 0.066667 for n = 8 is likewise a \emph{certified} optimum, the
total swing count having been swept over its full range to the
theoretical maximum T = \(8\cdot2^7\) = 1024. The exact target
\((2c,\dots,2c,c)\) is infeasible for 6 \(\le\) n \(\le\) 10 by the
computation above, so the \(L_1-optimum\) is strictly positive. It is
the monotone decrease \(n\cdot L_1\) = 0.682 \(\to\) 0.598 \(\to\) 0.533
that supports the \emph{conjectured} \(\Theta(1/n)\) asymptotics for all
n; deriving an analytic lower bound
\(\|\mathrm{Bz}(v)\ -\ \psi^n\|_1\ \ge\) c/n with an effective constant
--- not merely certification on the first three scales --- remains open.
The optimal profiles are asymmetric: the two lightest voters --- the
special voter and one further voter --- carry the reduced swing count,
and they do not sit at the ideal scale T = \((2n-1)c\) (n = 6: profile
(6,4,6,6,6,4) at weights (3,1,3,3,3,1), quota 11; n = 7: profile
(5,7,7,7,7,7,5) at weights (2,64,64,64,64,64,1), quota 65; n = 8:
profile (6,8,8,8,8,8,8,6) at weights (2,3,3,3,3,3,3,1), quota 4). This
is consistent with Corollary 6: in the exact case the special voter
alone would be uniquely lightest, whereas near-optimal games distribute
the reduction between the two lightest voters.

\subsection{10. Computational methodology: exact construction and
certified
obstruction}\label{computational-methodology-exact-construction-and-certified-obstruction}

This section collects the two computational engines behind the paper as
a reusable methodology --- with the complexity statements and instance
tables expected of an operations-research study. Every claim below is
backed by a named script and a full log in the companion code package
(Section 11).

\subsubsection{10.1 Constructive engine: closed forms and the W-family
MILP}\label{constructive-engine-closed-forms-and-the-w-family-milp}

\emph{Task.} Given n, produce an explicit simple game with Bz(v) =
\(\psi^n,\) or show that the family under consideration cannot cover n.

\begin{enumerate}
\def\labelenumi{\arabic{enumi}.}
\tightlist
\item
  \textbf{Residue arithmetic.} For the complement-of-edges family,
  realization reduces (Theorem 1) to the diophantine equation 7c =
  \(d(2n-8)\) + \(4(n-1)\) in positive integers d, c with 1 \(\le\) d
  \(\le\ n-2,\) d + c/2 \(\le\ n-1,\) together with Erdős--Gallai
  graphicality of the degree sequence \((d,\dots,d,\) d+c/2). This is an
  O(1) decision per n, and the family covers exactly the four residue
  classes n \(\equiv\) 0, 2, 3, 6 (mod 7) (Proposition 8).
\item
  \textbf{W-family MILP/CP-SAT.} For the remaining residue classes,
  realization over the W-family of Proposition 11 is formulated as a
  MILP with edge variables \(x_e\) (e \(\in\) G) and triple variables
  \(y_T\) (T \(\in\) F), the bilinear products z = \(y\cdot x\)
  linearized in the standard way, and the exact swing formula (U)
  together with the monotonicity constraint Pairs(F) \(\subseteq\) G as
  linear constraints. \textbf{Model size: \(O(n^3)\) variables and
  constraints} (C(n,2) + C(n,3) variables) --- versus \(2^n\) for the
  general monotone Boolean MILP of Section 7, which becomes intractable
  near n = 18. The compressed model solves every instance up to n = 64
  in seconds to minutes (instance data below); its practical limiting
  factor is memory, reached around n \(\approx\) 80.
\item
  \textbf{Closed forms.} Proposition 13 gives explicit triangle
  constructions for the three hard residue classes n \(\equiv\) 1, 4, 5
  (mod 7), valid for all n \(\ge\) 22, 18, 19 respectively; the small
  values 8, 10, 11, 12, 15 are handled by the \(\deg_G\) + \(r_i\) = D
  structure (Corollary 16). Combined with Step 1 and Corollary 14,
  \textbf{every n \(\ge\) 6 admits an explicit closed form} (Corollary
  17).
\end{enumerate}

Representative W-family instances (hard residue classes; c = scale
factor). All are explicit closed forms, re-verified by direct swing
enumeration against \(\psi^n\) (full data in appendix-data-n80.json):

{\def\LTcaptype{none} % do not increment counter
\begin{longtable}[]{@{}llllll@{}}
\toprule\noalign{}
n & n mod 7 & c & edges & triples & construction \\
\midrule\noalign{}
\endhead
\bottomrule\noalign{}
\endlastfoot
18 & 4 & 24 & 33 & 6 & deg+r = D = 4 \\
22 & 1 & 26 & 32 & 5 & tri(D3,f5,q3) \\
25 & 4 & 30 & 36 & 6 & tri(D3,f6,q3) \\
39 & 4 & 50 & 62 & 6 & tri(D3,f6,q2) \\
67 & 4 & 90 & 114 & 6 & tri(D3,f6,q0) \\
71 & 1 & 96 & 123 & 5 & tri(D3,f5,q0) \\
74 & 4 & 100 & 127 & 6 & tri(D3,f6,q0) \\
78 & 1 & 106 & 136 & 5 & tri(D3,f5,q0) \\
\end{longtable}
}

Here \texttt{tri(D,f,q)} denotes the triangle construction of
Proposition 13 with base degree D, f disjoint triangles and q special
triples (containing the special player), and \texttt{deg+r\ =\ D} the
small-value structure of Corollary 16. The structural insight for the
hard residues is that the solutions use few exceptional triples --- f
\(\in\ \{4,\) 5, \(6\}\) --- and that for n \(\ge\) 67 the special
player need not participate in any special triple (q = 0; Corollary 15).

\subsubsection{10.2 Verification engine: certification of weighted
obstructions}\label{verification-engine-certification-of-weighted-obstructions}

\emph{Task.} Given n, decide whether any weighted majority game {[}q;
w{]} on n voters realizes the exact target swing vector
\((2c,\dots,2c,c);\) if none exists, produce a certificate.

\begin{enumerate}
\def\labelenumi{\arabic{enumi}.}
\tightlist
\item
  \textbf{Constraint system and conclusive weight bound.} The model
  (\texttt{conj15\_\allowbreak feas.\allowbreak py}) has variables
  \(w_1\ \le\ \dots\ \le\ w_n\ \le\) W, q, and the swing counts
  \(\eta_i,\) together with the exact profile equations \((\eta_i\) = 2c
  for i \(\ne\) s, \(\eta_s\) = c) and the structural consequences (L),
  \((M')\) of Section 9. Infeasibility at W = \(2^{n-1}\) is conclusive,
  because every weighted majority game on n voters admits an integer
  representation with all weights in {[}1, \(2^{n-1}]\) (Muroga's
  sharper bound \(3^{n/3}\) applies for n \(\ge\) 3). A solver
  certificate of infeasibility at this bound is therefore a proof for
  that n, not a heuristic.
\item
  \textbf{Symmetry breaking.} Non-special voters may be relabeled in
  nondecreasing weight order (they carry equal target swing 2c). Without
  this constraint n = 8 remains inconclusive after 20 min; with it the
  certificate is found in 316 s.
\item
  \textbf{Mirror pruning.} The sharp constraints \((M')\) of Corollary 7
  are consequences of the exact hypothesis, hence valid; enforcing them
  (\texttt{conj15\_\allowbreak feas2.\allowbreak py}) yields a roughly
  six-fold speedup (n = 6: 2.5 s vs 2.9 s; n = 7: 19 s vs 117 s, a
  six-fold speedup) and --- decisively --- makes the n \(\ge\) 9
  frontier decidable at all.
\end{enumerate}

Certified obstructions (deterministic, at W = \(2^{n-1}):\)

{\def\LTcaptype{none} % do not increment counter
\begin{longtable}[]{@{}lllll@{}}
\toprule\noalign{}
n & W = \(2^{n-1}\) & model & verdict & runtime \\
\midrule\noalign{}
\endhead
\bottomrule\noalign{}
\endlastfoot
6 & 32 & standard & INFEASIBLE & 2.9 s \\
7 & 64 & standard & INFEASIBLE & 117 s \\
8 & 128 & +symmetry breaking & INFEASIBLE & 316 s \\
9 & 256 & +mirror & INFEASIBLE & 1562.9 s \\
10 & 512 & +mirror & INFEASIBLE & 9527.4 s \\
\end{longtable}
}

The n = 6 verdict agrees with exhaustive enumeration of all 1,111
weighted games (Tier 1); n = 7--10 are Tier-2 deterministic
certificates. Under a two-hour budget the standard (unpruned) model
remains UNKNOWN at n = 9 and n = 10 --- the mirror constraints are what
closes these cases. For n \(\ge\) 11 the same solver is run budgeted, so
a negative outcome there is heuristic (Tier 3) and is labelled as such.
The mirror constraints --- 3584 at n = 9 and 9216 at n = 10,
\(i.e. C(n-1,2)\cdot2^{n-2}\) window conditions, one per subset of the
\(n-3\) remaining voters and for both windows --- are themselves a
measure of the rigidity that \((M')\) forces on any realizing game.

\subsubsection{10.3 Exact optimization over total swing
count}\label{exact-optimization-over-total-swing-count}

For the quantitative form of Conjecture 15 we compute the minimal
\(L_1-distance\) \(\|\mathrm{Bz}(v)\ -\ \psi^n\|_1\) exactly. Fixing the
total number of swings T = \(\sum\eta_i,\) the integer objective
\(\sum_{i\ne s}\ \vert(2n-1)\eta_i\ -\ 2T\vert\) +
\(\vert(2n-1)\eta_s\ -\ T\vert\) is minimized in CP-SAT; T is enumerated
over its full range up to the theoretical maximum \(n\cdot2^{n-1},\) at
the same conclusive weight bound
(\texttt{conj15\_\allowbreak l1.\allowbreak py},
\texttt{conj15\_\allowbreak l1\_\allowbreak range.\allowbreak py}). The
method reproduces the exhaustive n = 6 optimum 5/44 (weights
(3,1,3,3,3,1), quota 11) exactly --- a strong cross-check of both
methods. Certified optima over the full swing range: n = 7,
\(L_1\ \approx\) 0.085470 at T = 45; n = 8, \(L_1\ \approx\) 0.066667 at
T = 60 (over T \(\le\) 1024). These give c \(\le\) min(30/44,
\(7\cdot0.08547,\ 8\cdot0.066667)\ \approx\) 0.533. At n = 9 the problem
hits the CP-SAT complexity cliff: with 120 s per T-value, 23 of the 112
values in {[}9,120{]} are decided, 0 are certified infeasible, and the
best found \(L_1\ \approx\) 0.1042 (at T = 70) is not certified. We
report this honestly as the current frontier: \textbf{c \(\le\) 0.533 is
certified for n \(\le\) 8; no certified value exists at n = 9.}

\subsubsection{10.4 Reproducibility}\label{reproducibility}

Every claim above is tied to a named script ---
\texttt{conj15\_\allowbreak feas.\allowbreak py},
\texttt{conj15\_\allowbreak feas2.\allowbreak py},
\texttt{conj15\_\allowbreak l1.\allowbreak py},
\texttt{conj15\_\allowbreak l1\_\allowbreak range.\allowbreak py},
\texttt{weighted\_\allowbreak enum.\allowbreak py},
\texttt{probe6\_\allowbreak identity.\allowbreak py}, and the W-family
MILP/CP-SAT solver of Proposition 12 --- with full logs in the companion
package; solver versions, hardware, and per-run parameters are recorded
in the package README.

\subsection{11. The W-family: unified constructions, structural
properties, and open
problems}\label{the-w-family-unified-constructions-structural-properties-and-open-problems}

This section develops the W-family into its full scope: it proves the
unified swing formula that the construction half (Section 8) relies on,
establishes the structural properties that govern which residue classes
admit closed-form realizations, and collects the open problems that
remain. All results are labelled with their derivation or verification
method; unproven numerical observations are explicitly marked as such.
Reproducible code is given in the appendix.

\subsubsection{Proposition 8 (Exact description of complement-of-edges
coverage; Erdős--Gallai
computation)}\label{proposition-8-exact-description-of-complement-of-edges-coverage-erdux151sgallai-computation}

For each n \(\ge\) 6, the complement-of-edges family (Theorem 1) is
solvable at n if and only if there are positive integers d, c with 7c =
\(d(2n-8)\) + \(4(n-1),\) 1 \(\le\) d \(\le\ n-2,\) c even, d + c/2
\(\le\ n-1,\ (n\cdot d\) + c/2) even, and the degree sequence
\((d,\dots,d,\) d+c/2) passes the Erdős--Gallai test and is
constructible by Havel--Hakimi.

\textbf{Lemma A (two-value Erdős--Gallai reduction).} Let 1 \(\le\) d
\textless{} D \(\le\ n-1.\) The degree sequence (d, \(\dots,\) d, D)
with \(n-1\) copies of d is graphical if and only if \((n-1)d\) + D is
even, D \(\le\ n-1,\) and D \(\le\ d(n-d).\)

\emph{Proof.} Sorting descending as (D, d, \(\dots,\) d), the
Erdős--Gallai inequalities are, for every k,
\(\sum_{i\le k}\ a_i\ \le\ k(k-1)\) + \(\sum_{i>k}\ \min(a_i,\) k). For
k = 1 this reads D \(\le\ n-1.\) For 2 \(\le\) k \(\le\) d every \(a_i\)
= d \(\ge\) k, so the RHS is \(k(k-1)\) + \((n-k)k\) = \(k(n-1);\) the
strongest such inequality is at k = 2, i.e.~D \(\le\ 2(n-1)\ -\) d,
which is implied by D \(\le\ n-1.\) For k \textgreater{} d the RHS is
\(k(k-1)\) + \((n-k)d,\) and the strongest case is k = d+1, reading D
\(\le\) d(d+1) + \((n-d-1)d\) = \(d(n-d).\) Hence, besides parity and D
\(\le\ n-1,\) the only content is D \(\le\ d(n-d).\) (Sufficiency is the
Erdős--Gallai theorem.) \(\blacksquare\)

\emph{Analytic conclusion.} Applying Lemma A to the four residue classes
with the constant d above (n \(\equiv\) 6, 0, 2, 3 (mod 7); D = d + c/2
= (4n+4)/7, (5n+7)/7, (3n+1)/7, (6n+10)/7), the conditions hold for
every admissible n: for n \(\equiv\) 6 (d = 2) both D \(\le\ n-1\) and D
\(\le\ 2(n-2)\) hold for all n \(\ge\) 6; for n \(\equiv\) 0 (d = 3) D
\(\le\ 3(n-3)\) always and D \(\le\ n-1\) for n \(\ge\) 7; for n
\(\equiv\) 2 (d = 1) both hold for all n \(\ge\) 9; for n \(\equiv\) 3
(d = 4) D \(\le\ 4(n-4)\) always and D \(\le\ n-1\) for n \(\ge\) 17.
Parity: \((n-1)d\) + D equals (18m+14), \((26m-2),\) (10m+2), (34m+12)
respectively for n = 7m+r, all even. Hence the complement-of-edges
family covers n \(\equiv\) 6 (mod 7) from n \(\ge\) 6, n \(\equiv\) 0
from n \(\ge\) 7, n \(\equiv\) 2 from n \(\ge\) 9, and n \(\equiv\) 3
from n \(\ge\) 17 --- with no computational verification. (The value n =
10, the only n \(\equiv\) 3 with 6 \(\le\) n \textless{} 17, is covered
by Corollary 16.)

\textbf{Corollary 9 (new infinite families).} All admissible n in the
following four residue classes are realized by complement-of-edges
constructions: n \(\equiv\) 6 (mod 7), d = 2; n \(\equiv\) 0 (mod 7), d
= 3; n \(\equiv\) 2 (mod 7), d = 1; n \(\equiv\) 3 (mod 7), d = 4. Each
construction is a complement-of-edges game whose graph is ``a star
centered at the special player plus a structured graph''; explicit
graphs for n = 20, 21, 23, 24 etc. were built and independently
re-checked by both the swing formula and (for small n) exhaustive
enumeration. The complement-of-edges family therefore covers \textbf{4/7
of all integers}, far beyond the two residue classes of the main text.

Uncovered residue classes: n \(\equiv\) 1, 4, 5 (mod 7) (n = 18, 19, 22,
25, 26, \(\dots\) in the main text).

\subsubsection{\texorpdfstring{Necessary property: all \((n-1)-sets\)
are
winning}{Necessary property: all (n-1)-sets are winning}}\label{necessary-property-all-n1-sets-are-winning}

\textbf{Proposition 10 (necessary; computational verification + proof
sketch).} If the Banzhaf index of a simple game v equals \(\psi^n,\)
then every \((n-1)-element\) coalition is winning.

\emph{Verification}: exhaustive enumeration of all 7,828,354 simple
games on 6 voters (all 360 solutions satisfy it); for n = 7, 8 the MILP
becomes infeasible when the constraint ``some \((n-1)-set\) is losing''
is added.

\emph{Proof sketch}: if \(N\setminus\{x\}\) is losing, player x swings
on S = \(N\setminus\{x\};\) moreover, on each \((n-2)-set\) S =
\(N\setminus\{x,j\},\allowbreak\) player x may swing (because
\(N\setminus\{j\}\) is winning) while player j never does (because
\(N\setminus\{x\}\) is losing), so x accumulates extra swings at the top
levels, contradicting \(\beta\ \propto\ (2,\dots,2,1).\) A complete
proof needs a level-by-level count; left as an unfinished proof in the
work record.

\emph{Use}: it shrinks the search space to games in which all
\((n-1)-sets\) win (adopted in Task F).

\subsubsection{A unified framework: the
W-family}\label{a-unified-framework-the-w-family}

\textbf{Proposition 11 (unifying theorem).} For a graph G and a family F
of triples with Pairs(F) \(\subseteq\) G (monotonicity), define
v(G,F)(S) = 1 \(\iff\ \vert S\vert\ \ge\ n-1,\) or \((\vert S\vert\) =
\(n-2\) and N\S \(\in\) G), or \((\vert S\vert\) = \(n-3\) and
N\S \(\in\) F). Then v(G,F) is a simple game and \(\beta_i(v(G,F))\) =
\(\vert G\vert\) + \((n-1)\) +
\(\vert F\vert\ -\ 2\cdot\deg_G(i)\ -\ 2\cdot r_i\) + \(e_i,\) (U) where
\(r_i\) = \(\vert\{X\ \in\) F : i \(\in\ X\}\vert\) and \(e_i\) =
\(\vert\{X\ \in\) F : i \(\in\) X,
\(X\setminus\{i\}\ \notin\ G\}\vert.\)

\begin{itemize}
\tightlist
\item
  With G = E and F = \(\varnothing,\) (U) reduces to Theorem 1
  (complement-of-edges);
\item
  With G = complement(E) and arbitrary F, (U) reduces to Theorem 2
  (threshold-with-exceptions).
\end{itemize}

\emph{Verification}: for n = 6..10, (U) agrees with direct swing counts
on random monotone (G,F) instances. This unification shows that the two
construction families belong to one parametric framework, and provides a
single equation for finding new constructions.

\subsubsection{\texorpdfstring{Proposition 12 (Task F: an
\(O(n^3)-variable\) MILP over the W-family; verification extended to n
\(\ge\)
18)}{Proposition 12 (Task F: an O(n\^{}3)-variable MILP over the W-family; verification extended to n \textbackslash ge 18)}}\label{proposition-12-task-f-an-onuxb3-variable-milp-over-the-w-family-verification-extended-to-n-18}

The general MILP of Section 7 has \(2^n\) variables and becomes
infeasible around n = 18. Using Proposition 11's W-family, we compress
the search to variables \(x_e\) (edge e \(\in\) G, C(n,2) of them) and
\(y_T\) (triple T \(\in\) F, C(n,3) of them), with the bilinear terms z
= \(y\cdot x\) linearized in the standard way, yielding a MILP with
\textbf{\(O(n^3)\) variables} (formula (U) and monotonicity Pairs(F)
\(\subseteq\) G as linear constraints). This MILP solves all instances
up to n \(\le\) 64 in seconds to minutes; every resulting construction
is an explicit closed form (Corollaries 14--17, Proposition 13),
re-verified by formula (U). The scales c are:

\(\{\)\def\LTcaptype{none} \% do not increment counter

\begin{longtable}[]{@{}llll@{}}
\toprule\noalign{}
n & n mod 7 & c & construction \\
\midrule\noalign{}
\endhead
\bottomrule\noalign{}
\endlastfoot
18 & 4 & 24 & deg+r=D \\
19 & 5 & 22 & tri(D3,f4,q2) \\
20 & 6 & 20 & comp-edge \\
21 & 0 & 26 & comp-edge \\
22 & 1 & 26 & tri(D3,f5,q3) \\
23 & 2 & 18 & comp-edge \\
24 & 3 & 36 & comp-edge \\
\end{longtable}

\(\}\)

\(\{\)\def\LTcaptype{none} \% do not increment counter

\begin{longtable}[]{@{}llll@{}}
\toprule\noalign{}
n & n mod 7 & c & construction \\
\midrule\noalign{}
\endhead
\bottomrule\noalign{}
\endlastfoot
25 & 4 & 30 & tri(D3,f6,q3) \\
26 & 5 & 32 & tri(D3,f4,q2) \\
27 & 6 & 28 & comp-edge \\
28 & 0 & 36 & comp-edge \\
29 & 1 & 36 & tri(D3,f5,q2) \\
30 & 2 & 24 & comp-edge \\
31 & 3 & 48 & comp-edge \\
\end{longtable}

\(\}\)

Here \texttt{tri(D,f,q)} and \texttt{comp-edge} are as in Proposition 13
/ Section 10.1.

Then n = 32 \((\equiv4)\) c=40, n=33 \((\equiv5)\) c=42, n=34
\((\equiv6)\) c=36, n=35 \((\equiv0)\) c=46, n=36 \((\equiv1)\) c=46;
n=37 \((\equiv2)\) c=30 and n=38 \((\equiv3)\) c=60 by
complement-of-edges; n=39 \((\equiv4)\) c=50. Further: \(40(\equiv5)\)
c=52, \(43(\equiv1)\) c=56, \(46(\equiv4)\) c=60, \(47(\equiv5)\) c=62,
\(50(\equiv1)\) c=66, \(53(\equiv4)\) c=70, \(54(\equiv5)\) c=72,
\(57(\equiv1)\) c=76, \(60(\equiv4)\) c=80, \(61(\equiv5)\) c=82,
\(64(\equiv1)\) c=86, all rechecked by (U).

\textbf{Closed forms cover every n \(\ge\) 67 (supplement to Proposition
12).} Corollary 15 shows that for n \(\ge\) 67 the special player need
not participate in any F-triple (q = 0). Applying the triangle
construction directly yields the remaining values: n = 67 \((\equiv4)\)
c = 90, n = 68 \((\equiv5)\) c = 92, n = 71 \((\equiv1)\) c = 96, n = 74
\((\equiv4)\) c = 100, n = 75 \((\equiv5)\) c = 102, n = 78
\((\equiv1)\) c = 106, with \(\vert G\vert\) = 123, \(\vert F\vert\) = 5
at n = 71 and \(\vert G\vert\) = 136, \(\vert F\vert\) = 5 at n = 78.
\textbf{Hence every n with 6 \(\le\) n \(\le\) 80 is covered without
gaps, for all residue classes mod 7} (verified: 75/75 values in
{[}6,80{]}). The structural insight for the hard residues is that their
solutions use few exceptional triples: the triangle construction needs
at most six \((\vert F\vert\ \in\ \{4,\) 5, \(6\}),\) and for n \(\ge\)
67 none of them involves the special player. For n \(\ge\) 81 the
\(O(n^3)\) model exhausts memory around n \(\approx\) 80, recorded as
the second computational lower bound of this method.

\emph{Methodological remark}: the speed of the W-family MILP/CP-SAT (n =
18 in about 5 s) shows that the unified framework of Proposition 11 not
only unifies the constructions but also gives a scalable computational
route; this is a direct payoff of Tasks C/E.

\subsubsection{Closed-form constructions for the hard residue classes;
full
resolution}\label{closed-form-constructions-for-the-hard-residue-classes-full-resolution}

\textbf{Proposition 13 (found by exploration, verified via the proven
formula (U)).} For each of the three residue classes n \(\equiv\) 1, 4,
5 (mod 7), there is an explicit closed-form construction realizing
\(\psi^n,\) valid for all n from a small threshold:

{\def\LTcaptype{none} % do not increment counter
\begin{longtable}[]{@{}
  >{\raggedright\arraybackslash}p{(\linewidth - 8\tabcolsep) * \real{0.1311}}
  >{\raggedright\arraybackslash}p{(\linewidth - 8\tabcolsep) * \real{0.0656}}
  >{\raggedright\arraybackslash}p{(\linewidth - 8\tabcolsep) * \real{0.1967}}
  >{\raggedright\arraybackslash}p{(\linewidth - 8\tabcolsep) * \real{0.2787}}
  >{\raggedright\arraybackslash}p{(\linewidth - 8\tabcolsep) * \real{0.3279}}@{}}
\toprule\noalign{}
\begin{minipage}[b]{\linewidth}\raggedright
residue
\end{minipage} & \begin{minipage}[b]{\linewidth}\raggedright
f
\end{minipage} & \begin{minipage}[b]{\linewidth}\raggedright
c
\end{minipage} & \begin{minipage}[b]{\linewidth}\raggedright
special degree s
\end{minipage} & \begin{minipage}[b]{\linewidth}\raggedright
valid for
\end{minipage} \\
\midrule\noalign{}
\endhead
\bottomrule\noalign{}
\endlastfoot
n \(\equiv\) 1 & 5 & \((10(n-1)-28)/7\) & (c+6)/2 \(-\) q & n \(\ge\)
22 \\
n \(\equiv\) 4 & 6 & \((10(n-1)-30)/7\) for n \(\ge\) 25 (D = 3); c = 24
at n = 18 (D = 4) & D + c/2 \(-\) q & n \(\ge\) 18 (D = 4 at n = 18,
else D = 3) \\
n \(\equiv\) 5 & 4 & \((10(n-1)-26)/7\) & (c+6)/2 \(-\) q & n \(\ge\)
19 \\
\end{longtable}
}

Here q \(\in\ \{0,\dots,f\}\) is the number of \emph{special triples}
--- F-triples containing the special player (equal to the number of
triples in which the special player participates). The construction
takes f disjoint triangles: q of them of the form \(\{n,\) a, \(b\}\)
(special triple) and p = f \(-\) q of them involving only non-special
players; every player in a triangle has degree 2 in G (base degree D =
3, so each triangle player has degree \(D-1\) = 2) and lies in exactly
one F-triple \((r_i\) = 1, \(e_i\) = 0). The remaining \(n-1-3f+q\)
non-special players have degree D = 3 and \(r_i\) = 0; the special
player has degree s = D + c/2 \(-\) q (for D = 3 this is (c+6)/2 \(-\)
q). By formula (U) the swing counts are \(\beta_i\) = \(\vert G\vert\) +
\((n-1)\) + f \(-\) 6 = 2c (non-special), \(\beta_n\) = \(\vert G\vert\)
+ \((n-1)\) + f \(-\) 2s \(-\) 2q = c. The smallest feasible q admits
the closed form q = max(0, \(\lceil(K\ -\ n)/14\rceil),\) K = 57 (n
\(\equiv\) 1), 67 (n \(\equiv\) 4), 47 (n \(\equiv\) 5). (6) Indeed,
writing \(\delta\) = 2, 1, 3 for n \(\equiv\) 1, 4, 5 respectively,
direct substitution (D = 3, c = \((10(n-1)-A)/7\) with A = 28, 30, 26)
gives the special degree s = \((5n+\delta)/7\ -\) q and the residual
degree sequence \((3^a,\ 2^x)\) on the m = \(n-1-3f+q\) non-triangle
players, where x = s \(-\) 2q and a = m \(-\) x = \((2(n-K)\) + 28q)/7.
Since n \(\equiv\) K (mod 7), a is \textbf{always even}. The feasibility
conditions x \(\ge\) 0, a \(\ge\) 0, a \(\le\) (m+1)/2, q \(\le\) f, m
\(\ge\) 4, s \(\le\ n-1\) all hold at (6): x \(\ge\) 0 \(\iff\) q
\(\le\ (5n+\delta)/21\) and a \(\le\) (m+1)/2 \(\iff\) 49q \(\le\) 3n +
4K \(-\) 21f are the binding ones, each holding for every n \(\ge\ n_0\)
once the elementary inequalities are checked (e.g.~n
\(\ge\ (3K+42-2\delta)/13\) \textless{} \(n_0\) for x \(\ge\) 0; n
\(\ge\ (3.5K+49-(4K-21f))/6.5\) \textless{} \(n_0\) for the chord
condition); a \(\ge\) 0 is built into (6), q \(\le\) f and s
\(\le\ n-1\) are automatic. The residual graph \((3^a,\ 2^x)\) is then
constructed explicitly: place the a degree-3 vertices at the even
positions 0, 2, \(\dots,\ 2a-2,\) add the cycle (0 1 \(\dots\ m-1\) 0),
and add the matching chords (2j, 2j+2) for j = 0, 2, \(\dots,\ a-2.\)
Each even-position vertex is the endpoint of exactly one chord, so it
has degree 3; all other vertices have degree 2. (a \(\le\) (m+1)/2
\(\iff\ 2a-2\ \le\ m-1\) makes the positions exist; a even closes the
chord set.) Thus (6) realizes \(\psi^n\) for every n \(\equiv\) 1, 4, 5
(mod 7) with n \(\ge\) 22, 25, 19 respectively, with no search and no
reliance on numerical verification.

\textbf{Corollary 14 (every n \(\ge\) 18 has a closed form).} Since n
\(\equiv\) 0, 2, 3, 6 (mod 7) are covered by the complement-of-edges
construction of Proposition 8 (by Lemma A, analytically for every
admissible n) and n \(\equiv\) 1, 4, 5 by the general triangle
construction of Proposition 13 (by the closed form (6), analytically for
n \(\ge\) 22, 25, 19 respectively), \textbf{every n \(\ge\) 18 admits an
explicit closed-form simple game with Bz(v) = \(\psi^n\)}. For n = 18
this uses base degree D = 4 with six disjoint triangles, three of which
contain the special player: c = 24, \(\vert G\vert\) = 33,
\(\vert F\vert\) = 6.

\textbf{Corollary 15 (unified parametric family).} With base degree d =
3 for all non-special players and f = the smallest non-negative integer
with f \(\equiv\) 5(n mod 7) (mod 7), the same construction realizes
\(\psi^n\) for every residue class; for n \(\ge\) 67 the special player
need not participate in any F-triple, and for smaller n a suitable
number q of special triples is used. Thus a single parametric family
(the W-family of Proposition 11 with the degree/f-triple data above)
covers all n.

\textbf{Corollary 16 (closed forms for the five remaining small
values).} For n = 8, 10, 11, 12, 15 explicit closed-form constructions
are also found (via the same \(\deg_G\) + \(r_i\) = D structure, with
the data below; verified by formula (U)):

{\def\LTcaptype{none} % do not increment counter
\begin{longtable}[]{@{}llllll@{}}
\toprule\noalign{}
n & n mod 7 & c & \(\vert G\vert\) & \(\vert F\vert\) & D = deg+r
(non-special) \\
\midrule\noalign{}
\endhead
\bottomrule\noalign{}
\endlastfoot
8 & 1 & 8 & 15 & 2 & 4 \\
10 & 3 & 10 & 16 & 1 & 3 \\
11 & 4 & 14 & 22 & 6 & 5 \\
12 & 5 & 16 & 25 & 6 & 5 \\
15 & 1 & 20 & 32 & 2 & 4 \\
\end{longtable}
}

Each is a W-family game v(G,F) with all swing counts
\(\propto\ (2,\dots,2,1)\) and Pairs(F) \(\subseteq\) G (monotone), so
Bz(v) = \(\psi^n.\)

\textbf{Corollary 17 (full resolution --- analytic).} Combining
Corollary 14 (analytic, via Lemma A and (6)) with Corollary 16,
\textbf{every n \(\ge\) 6 admits an explicit closed-form simple game
with Bz(v) = \(\psi^n.\)} The proof is fully analytic: for n \(\equiv\)
0, 2, 3, 6 (mod 7) the complement-of-edges construction of Proposition 8
realizes \(\psi^n\) for every admissible n (Lemma A: the two-value
Erdős--Gallai conditions hold for all n in the residue class), and for n
\(\equiv\) 1, 4, 5 (mod 7) the triangle construction of Proposition 13
with the closed form (6) realizes \(\psi^n\) for every n \(\ge\) 22, 25,
19 respectively; the finitely many remaining values 6 \(\le\) n \(\le\)
21 are the explicit games of Corollary 16 (n = 8, 10, 11, 12, 15) and of
Proposition 8 (n = 6, 7, 9, 13, 14, 16, 17), with n = 18 the D = 4
variant of Corollary 14. Kurz's Conjecture 16 is therefore fully
resolved with explicit constructions for all n, with no reliance on
computer search for existence.

\subsubsection{Verification of the core theorems (independent
re-check)}\label{verification-of-the-core-theorems-independent-re-check}

The three central claims were re-verified independently of their
original derivations:

\textbf{(a) The W-family swing formula (U), Proposition 11.} The
identity \(\beta_i(v(G,F))\) = \(\vert G\vert\) + \((n-1)\) +
\(\vert F\vert\ -\ 2\cdot\deg_G(i)\ -\ 2\cdot r_i\) + \(e_i\) was
checked against direct swing enumeration on 440 randomly generated
monotone instances (Pairs(F) \(\subseteq\) G, n = 4..14): \textbf{0
mismatches}. The two special cases were checked separately --- Theorem 1
(F = \(\varnothing,\) 199 random instances) and Theorem 2 (G =
complement(E), 29 instances) --- also with \textbf{0 mismatches}.

\textbf{(b) The closed-form constructions of Proposition 13 (analytic by
Lemma A and (6)).} As an independent confirmation of the analytic proof,
the triangle-with-special-triples construction was generated and its
swing counts verified via (U) for every n \(\in\) {[}6, 50{]} with n
\(\equiv\) 1, 4, 5 (mod 7) and n \(\ge\) 18, and for the
complement-of-edges constructions for the other residue classes:
\textbf{all verified} (the five values 8, 10, 11, 12, 15 are covered by
the \(\deg_G\) + \(r_i\) = D constructions of Corollary 16, also
verified). The Erdős--Gallai conditions of Lemma A and the feasibility
inequalities of (6) were additionally checked by direct computation for
all n up to \(10^6\) --- agreeing with the analytic conclusion.

\textbf{(c) The necessary property of Proposition 10.} Exhaustive
enumeration for n = 6 (all 360 solutions) and infeasibility of the MILP
with a ``losing \((n-1)-set\)'' constraint for n = 7, 8 confirm that all
\((n-1)-sets\) are winning in any game with Bz(v) = \(\psi^n.\)

All verification scripts are included in the reproducible code package.

\subsubsection{Current coverage table}\label{current-coverage-table}

\(\{\)\def\LTcaptype{none} \% do not increment counter

\begin{longtable}[]{@{}lll@{}}
\toprule\noalign{}
n & method & c \\
\midrule\noalign{}
\endhead
\bottomrule\noalign{}
\endlastfoot
6 & compl & 4 \\
7 & compl & 6 \\
8 & deg+r & 8 \\
9 & compl & 6 \\
10 & deg+r & 10 \\
11 & deg+r & 14 \\
12 & deg+r & 16 \\
13 & compl & 12 \\
14 & compl & 16 \\
15 & deg+r & 20 \\
16 & compl & 12 \\
17 & compl & 24 \\
\end{longtable}

\(\}\)

(18--75 see the tables in Proposition 12.) n \(\equiv\) 0, 2, 3, 6 (mod
7) are covered for infinitely many n by Proposition 8; n \(\equiv\) 1,
4, 5 (mod 7) are covered by the triangle construction of Proposition 13
and by Corollary 16 for the smallest values. \textbf{Every n \(\ge\) 6
has a closed-form construction.} Full explicit (G, F) data for 6 \(\le\)
n \(\le\) 80 are provided in machine-readable form in
appendix-data-\texttt{n80.\allowbreak json}; entries for 41 \(\le\) n
\(\le\) 80 are produced by the closed-form constructions of Proposition
8 (n \(\equiv\) 0, 2, 3, 6 (mod 7)) and Proposition 13 (n \(\equiv\) 1,
4, 5 (mod 7)), and every entry (6 \(\le\) n \(\le\) 80) is re-verified
by direct enumeration of the swing counts --- equal to \(\psi^n,\) i.e.
\((2c,\dots,2c,c)\) --- and of the monotonicity condition Pairs(F)
\(\subseteq\) G.

\subsubsection{Open problems}\label{open-problems}

\begin{enumerate}
\def\labelenumi{\arabic{enumi}.}
\tightlist
\item
  Can Proposition 10 (all \((n-1)-sets\) winning) be proved rigorously?
\item
  Is c necessarily even?
\item
  What is the minimal c for each n? (Certified by the CP-SAT model of
  Section 7: n=6: 4, n=7: 6, n=8: 8, n=9: 6, n=10: 10.) For n = 8 the
  previously open ``7 or 8'' gap is closed: the infeasibility of every
  smaller even scale c \(\in\ \{2,\) 4, \(6\}\) is certified, and c = 8
  is realized (Table 2); likewise n = 9 (c \(\in\ \{2,\ 4\}\)
  infeasible, c = 6 realized) and n = 10 (c \(\in\ \{2,\) 4, 6, \(8\}\)
  infeasible, c = 10 realized). Note that c is even (Corollary 5), so
  the odd candidate c = 7 is already excluded; the certified values are
  conclusive because CP-SAT proves infeasibility exactly.
\item
  Is there a single uniform construction that realizes \(\psi^n\) for
  all n with no case distinction? The current resolution uses
  residue-specific parameters within one framework (Corollary 15).
\item
  \emph{(Weighted case, Conjecture 15.)} Can the certified infeasibility
  of the exact target \((2c,\dots,2c,c)\) for 6 \(\le\) n \(\le\) 10
  (Section 9) be turned into an analytic proof for all n? In particular,
  the quantitative form asks for a universal c \textgreater{} 0 with
  \(\|\mathrm{Bz}(v)\ -\ \psi^n\|_1\ \ge\) c/n for every weighted
  majority game v. The certified optimal distances 5/44, 0.085470,
  0.066667 (n = 6, 7, 8) give the empirical constant c \(\le\) 0.533; an
  analytic argument with c = \(\Theta(1)\) would settle the conjecture.
\end{enumerate}

\subsection{12. Application: an exact computational case
study}\label{application-an-exact-computational-case-study}

We close with a concrete case study applying the two engines of Sections
10.1--10.3 to a real voting body: the weighted voting rules of the
Council of the European Union (EU) and its predecessors. Every claim
below is tagged with its epistemic status --- \textbf{certified}
(deterministic computation at a conclusive weight bound), \textbf{exact}
(closed form / subset enumeration), or \textbf{contextual}
(institutional background, not a theorem of this paper).

\subsubsection{12.1 Actual Banzhaf power of historical EU Council
bodies}\label{actual-banzhaf-power-of-historical-eu-council-bodies}

The original European Economic Community (EEC) Council (Treaty of Rome,
1958) and its first enlargements were weighted majority games on n = 6,
9, 10 voters --- exactly the certified range of Section 9. Their
official qualified-majority weights and quotas (Felsenthal \& Machover
1998, who tabulate the EU Council history; the accession treaties) and
the exact Banzhaf power distributions are:

{\def\LTcaptype{none} % do not increment counter
\begin{longtable}[]{@{}
  >{\raggedright\arraybackslash}p{(\linewidth - 8\tabcolsep) * \real{0.2000}}
  >{\raggedright\arraybackslash}p{(\linewidth - 8\tabcolsep) * \real{0.2000}}
  >{\raggedright\arraybackslash}p{(\linewidth - 8\tabcolsep) * \real{0.2000}}
  >{\raggedright\arraybackslash}p{(\linewidth - 8\tabcolsep) * \real{0.2000}}
  >{\raggedright\arraybackslash}p{(\linewidth - 8\tabcolsep) * \real{0.2000}}@{}}
\toprule\noalign{}
\begin{minipage}[b]{\linewidth}\raggedright
body
\end{minipage} & \begin{minipage}[b]{\linewidth}\raggedright
n
\end{minipage} & \begin{minipage}[b]{\linewidth}\raggedright
weights w
\end{minipage} & \begin{minipage}[b]{\linewidth}\raggedright
quota q
\end{minipage} & \begin{minipage}[b]{\linewidth}\raggedright
\(L_1\) to \(\psi^n\)
\end{minipage} \\
\midrule\noalign{}
\endhead
\bottomrule\noalign{}
\endlastfoot
EEC-6 (1958--1973) & 6 & \((4\times3,\ 2\times2,\) 1) & 12 & 26/77
\(\approx\) 0.338 \\
EEC-9 (1973--1981) & 9 & \((10\times4,\ 5\times2,\ 3\times2,\) 2) & 41 &
2136/5389 \(\approx\) 0.396 \\
EEC-10 (1981--1986) & 10 & \((10\times4,\ 5\times3,\ 3\times2,\) 2) & 45
& 2528/6023 \(\approx\) 0.420 \\
EEC-12 (1986--1995) † & 12 & \((10\times4,\) 8, \(5\times4,\ 3\times2,\)
2) & 54 & 878/2323 \(\approx\) 0.378 \\
EU-15 (1995--2004) † & 15 & \((10\times4,\) 8,
\(5\times4,\ 4\times3,\ 3\times2,\) 2) & 62 & 98402/255867 \(\approx\)
0.385 \\
\end{longtable}
}

\emph{Weights are in compact multiplicity form: \(e.g. 10\times4\) means
four members of weight 10, a bare value means multiplicity one. The
exact Banzhaf swing-count vectors \(\beta\) and the derived power
distributions (percentages of total swings), in the same notation, are:
EEC-6 \(\beta\) = \((10\times3,\ 6\times2,\) 0), power =
\((23.8\times3,\ 14.3\times2,\) 0)\%; EEC-9 \(\beta\) =
\((53\times4,\ 29\times2,\ 21\times2,\) 5), power =
\((16.7\times4,\ 9.2\times2,\ 6.6\times2,\) 1.6)\%; EEC-10 \(\beta\) =
\((100\times4,\ 52\times3,\ 26\times3),\) power =
\((15.8\times4,\ 8.2\times3,\ 4.1\times3)\%;\) EEC-12 \(\beta\) =
\((286\times4,\) 242, \(148\times4,\ 102\times2,\) 40), power =
\((12.9\times4,\) 10.9, \(6.7\times4,\ 4.6\times2,\) 1.8)\%; EU-15
\(\beta\) = \((1968\times4,\) 1606,
\(1028\times4,\ 806\times3,\ 616\times2,\) 406), power =
\((11.2\times4,\) 9.1, \(5.8\times4,\ 4.6\times3,\ 3.5\times2,\)
2.3)\%.}

† \emph{Beyond the certified range} (n \(\ge\) 11): exact infeasibility
is not certified for these bodies (Tier-3 runs are heuristic, Section
9); the power distributions and \(L_1\) values are exact computations
reported as scaling context only, carrying no certificate.

Here \(L_1\) is the distance to \(\psi^n\) minimized over which member
occupies the half-power position of the egalitarian-with-one-exception
target. Three observations:

\begin{itemize}
\tightlist
\item
  \textbf{EEC-6 has a null player.} Luxembourg (weight 1, quota 12 of
  total 17) has zero swings: no coalition of the remaining five members
  sums to weight 11, so Luxembourg's vote is never decisive. Its power
  is 0, while the three large members (Germany, France, Italy) jointly
  hold 71.4\% of all power. This is a known pathology of the original
  weighted Council, reproduced here exactly.
\item
  \textbf{Real power is far from egalitarian.} The three bodies sit at
  \(L_1-distance\) 0.338--0.420 from \(\psi^n,\) and the power spread
  (largest/smallest positive power) ranges from \(3.85\times\) (EEC-10)
  to \(10.6\times\) (EEC-9).
\item
  \textbf{All three lie in the certified range} (n = 6, 9, 10), so the
  obstruction of Section 9 applies to them directly (exact
  infeasibility; the distance values above are exact subset-enumeration
  computations, reproduced by
  \texttt{case\_\allowbreak study\_\allowbreak eu.\allowbreak py} in the
  companion package).
\end{itemize}

\subsubsection{12.2 What the certified machinery tells a
designer}\label{what-the-certified-machinery-tells-a-designer}

For each body with n \(\le\) 10, Section 9 certifies ---
deterministically, at the conclusive weight bound \(w_i\ \le\ 2^{n-1}\)
--- that \textbf{no weighted majority game on n voters realizes
\(\psi^n\) exactly}. The decision-support consequences for an
institutional designer are:

\begin{enumerate}
\def\labelenumi{\arabic{enumi}.}
\tightlist
\item
  \textbf{Certified infeasibility of the target.} No re-weighting of
  these councils --- indeed no weighted rule on \(\le\) 10 voters at all
  --- can attain the egalitarian target \(\psi^n.\) A designer who
  insists on a weighted rule must accept a strictly positive distance to
  \(\psi^n.\)
\item
  \textbf{Certified best achievable (n \(\le\) 8).} Section 10.3
  computes the exact minimal \(L_1-distance\) over all weighted games:
  5/44 \(\approx\) 0.114 for n = 6 (attained at weights (3,1,3,3,3,1),
  quota 11, and verified against exhaustive enumeration of all 1,111
  weighted games), 0.0855 for n = 7, and 0.0667 for n = 8, each
  certified as a global optimum over the full total-swing range. The
  actual EEC-6 rule sits at \(L_1\) = 26/77 \(\approx\) 0.338 --- so
  even within the weighted class, a redesign could improve the distance
  to \(\psi^6\) by a factor of about three.
\item
  \textbf{Boundary of certification.} For n \(\ge\) 11 the exact
  infeasibility is not certified (Tier 3 runs are heuristic; Section 9).
  The later EEC-12 (n = 12) and EU-15 (n = 15) bodies lie beyond the
  certified range; their exact power distributions in Section 12.1 are
  reported as scaling context only, and their distance values carry no
  certificate.
\end{enumerate}

\subsubsection{12.3 The design decision and institutional
context}\label{the-design-decision-and-institutional-context}

The obstruction of Section 9 is a structural statement about the
weighted class: egalitarian power is unattainable for n \(\le\) 10, and
the real EU bodies are far from it. This is consistent with the actual
trajectory of the institution --- \textbf{contextual} --- the Council
moved away from a single weighted quota in the Lisbon Treaty, adopting a
double-majority rule (at least 55\% of member states representing at
least 65\% of the population) that cannot be represented as one weighted
majority game. The certified obstruction supplies the formal
decision-support complement to that historical choice: it quantifies how
far weighted rules can be from an egalitarian target, certifies that
they cannot reach it for small bodies, and --- via Section 10.3 ---
returns the exact best weighted approximation when the designer insists
on staying in the weighted class.

\subsubsection{12.4 A prescriptive design-analysis
procedure}\label{a-prescriptive-design-analysis-procedure}

For a designer facing a target \(\psi\) on an n-voter body, the two
engines of Section 10 compose into a decision-support procedure:

\begin{itemize}
\tightlist
\item
  if n \(\le\) 10 and \(\psi\) = \(\psi^n,\) the Section 9 obstruction
  certifies the decision (infeasible in the weighted class; feasible by
  explicit construction in the general class, Section 5);
\item
  if n \(\le\) 8 and \(\psi\) = \(\psi^n,\) the Section 10.3 optimizer
  returns the certified best weighted approximation and its distance;
\item
  otherwise the same solvers run under a time budget, and the outcome is
  labelled heuristic (Tier 3), not a certificate.
\end{itemize}

All numbers in this section are exact computations reproduced by
\texttt{case\_\allowbreak study\_\allowbreak eu.\allowbreak py} in the
companion code package.

\subsection{13. Conclusion}\label{conclusion}

We complement the construction half with a negative result in the
\emph{weighted} universe (Section 9): the egalitarian target \(\psi^n\)
cannot be reproduced by any weighted majority game. The obstruction is
anchored in a classical swap/block identity (Proposition 4) for
swing-count differences in weighted majority games, which forces the
distinguished player to carry minimal weight and concentrates every
other player's deviation into a rigid subset-sum window; the sharp
mirror form \((M')\) of those windows --- a new rigidity consequence of
the classical swap/block identity --- is our central structural
contribution. Deterministic CP-SAT, at the conclusive weight bound
\(w_i\ \le\ 2^{n-1},\) certifies infeasibility for n = 6, 7, 8, 9 and 10
(n = 6 consistent with the exhaustive enumeration of all 1,111 weighted
games, whose optimal \(L_1-distance\) to \(\psi^6\) is 5/44). We stress
the distinction between proof and certification: the infeasibility of
the exact target is a certified computation for 6 \(\le\) n \(\le\) 10,
not an analytic proof for all n, and likewise the quantitative form of
Conjecture 15 --- the universal bound
\(\|\mathrm{Bz}(v)\ -\ \psi^n\|_1\ \ge\) c/n --- is certified for n = 6,
7, 8 (optimal \(L_1\) distances 5/44, 0.085470, 0.066667, giving c
\(\le\) 0.533), while the \(\Theta(1/n)\) asymptotics for all n remains
a conjecture supported by these certified values. The two-tier
generalization \(\psi^n(3,1)\) is shown to be realizable by a
\emph{uniform} construction for all n \(\ge\) 6, while \(\psi^n(k,1)\)
is obstructed for k \(\ge\) 4. This gives a complete picture of the two
Kurz conjectures: Conjecture 16 holds by explicit construction, and
Conjecture 15 holds \emph{computationally} for n \(\le\) 10, with the
structural mechanism \((M')\) explaining why --- an analytic proof for
all n remains open.

As an operations-research contribution, the paper offers a reusable,
two-sided methodology rather than a one-off verification: a construction
engine (Section 10.1) that is closed-form for every n \(\ge\) 6 and
computationally verified for every 6 \(\le\) n \(\le\) 80, and an exact
verification engine (Section 10.2) that certifies non-existence at a
conclusive weight bound, with the structural property \((M')\) as its
pruning engine. Both halves are fully reproducible from the companion
code package, and both are explicit about their limits: the
constructions cover every n \(\ge\) 6, while the weighted obstruction is
certified for n \(\le\) 10 and the quantitative constant c \(\le\) 0.533
is certified for n \(\le\) 8 only. The Section 12 case study grounds
this methodology in a concrete institution --- the weighted Council of
the European Union --- turning the abstract inverse-optimization problem
into prescriptive decision support for the design of voting rules.

\subsection{Appendix. Additional explorations and negative
results}\label{appendix.-additional-explorations-and-negative-results}

The material in this appendix is kept for reproducibility and
completeness; none of it is needed for the main results. It documents
numerical observations and constructions that were explored and either
remain open or did not realize the target.

\textbf{Observation 7 (numerical).} In all known solutions the scaling
factor c is even (c \(\in\ \{4,\) 6, 8, 10, 12, 14, 16, 22, 24,
\(\dots\}).\) Whether this is forced is open. (In the weighted case
Corollary 5 resolves it: c must be even in any weighted realization.)

\textbf{Observation 9 (hidden structure of MILP solutions).} The MILP
solutions for n = 8, 10, 11, 12, 15 (Section 7) all belong to the
W-family (all \((n-1)-sets\) win). Their threshold parameters (G, F)
share: the special player has low degree in G (1--3) and appears in most
F-triples; the graph G is connected. These reflect the solver's specific
choices and show no complete symmetry/cyclic structure, indicating that
these n are ``essentially mixed'' and indeed cannot be covered by the
complement-of-edges family (Proposition 8) or any single uniform family
--- consistent with the arithmetic description of Proposition 8.
(Superseded for n \(\equiv\) 1, 4, 5 (mod 7) --- including n = 8, 11,
12, 15 above --- by the parametric triangle family of Corollary 15,
which covers those residue classes at every n by varying q; n = 10 (n
\(\equiv\) 3 (mod 7)) remains outside that family.)

\textbf{Unsuccessful constructions and negative results.}

\begin{enumerate}
\def\labelenumi{\arabic{enumi}.}
\tightlist
\item
  \textbf{``Core-set'' construction}: v(S)=1
  \(\iff\ \vert S\vert\ge n-1\) or \((\vert S\vert\ge n-3\) and
  \(C\subseteq S)\) for a fixed core C. For n = 6..10 no
  \(\vert C\vert\) yields \(\propto(2,\dots,2,1).\) Reason: the core
  makes swings too symmetric.
\item
  \textbf{Random sampling of the threshold (E,F) family}: for n = 8
  (known solvable), \(4\times10^6\) monotone random (E,F) samples gave
  zero hits --- the density of threshold-family solutions is essentially
  zero, so random search is hopeless and MILP is necessary. Recorded as
  a search lower bound.
\item
  \textbf{F = all triples containing the special player}: makes the
  special player a null player \((\beta_n\) = 0), so it does not realize
  \(\psi^n.\)
\item
  \textbf{E = \(K_{n-1}\) minus a regular graph}: makes
  \(\beta_n\ \equiv\) 0, so it does not realize \(\psi^n.\)
\item
  \textbf{Complement-of-edges at n = 18}: n \(\equiv\) 4 (mod 7), no
  integer (d,c) by Proposition 8. Confirmed.
\item
  \textbf{Closed-form parametric family for n \(\equiv\) 1, 4, 5 (mod
  7)}: CP-SAT was used to extract the full (G, F) solutions for n = 22,
  25, 26, 43. These solutions are highly ad-hoc: the special player's
  E-degree ranges from low (small n) to almost \(n-2\) (n = 43: degree
  42), \(\vert F\vert\) ranges from 4 to 299, and the c values are not
  unique across solvers (n = 43: HiGHS finds c = 56, CP-SAT c = 362,
  since c may be any multiple). \textbf{No clean parametric family could
  be extracted from these solutions}; this explains why these residue
  classes resist closed-form treatment (negative result, recorded).
  (Superseded: for n \(\equiv\) 1, 4, 5 (mod 7), Proposition 13 and
  Corollary 15 now derive the exact closed-form family from the
  deterministic certificates, so this negative conclusion no longer
  holds; the solver-dependence of the extracted c values observed here
  (n = 43: c = 56 vs 362) is explained there by the fact that c may be
  any multiple of the minimal feasible value.)
\end{enumerate}

\subsection{Declarations}\label{declarations}

\textbf{Declaration of competing interest.} The authors declare that
they have no known competing financial interests or personal
relationships that could have appeared to influence the work reported in
this paper.

\textbf{Funding.} This work was supported by the National Social Science
Fund of China under Grant No. 25CSH113.

\textbf{Data availability.} All results reported in this paper are
reproducible from the companion code package (CP-SAT/HiGHS models,
instance tables, and recorded runtimes), publicly available at

https://github.com/chenyongjingg/banzhaf-inverse-power

\textbf{CRediT authorship contribution statement.} Yongjin Chen:
Conceptualization, Methodology, Investigation, Writing --- original
draft. Shi Jin: Formal analysis, Software, Data curation, Writing ---
review \& editing. Songze Zhu: Conceptualization, Supervision, Writing
--- review \& editing, Funding acquisition.

\textbf{Declaration of generative AI and AI-assisted technologies in the
writing process.} During the preparation of this work the authors used
Claude (Anthropic) as an AI writing assistant to assist with language
editing, typesetting, and formatting of the manuscript. After using this
tool, the authors reviewed and edited the content as needed and take
full responsibility for the content of the publication.

\subsection{References}\label{references}

\begin{itemize}
\tightlist
\item
  Alon, N., \& Edelman, P. H. (2010). The inverse Banzhaf problem.
  \emph{Social Choice and Welfare}, 34(3), 371--377.
\item
  Chow, C. K. (1961). On the characterization of threshold functions. In
  \emph{Proceedings of the 2nd Annual Symposium on Switching Circuit
  Theory and Logical Design (SWCT 1961)} (pp.~34--38). IEEE.
  \url{https://doi.org/10.1109/FOCS.1961.24}
\item
  de Keijzer, B., Klos, T., \& Zhang, Y. (2014). Finding optimal
  solutions for voting game design problems. \emph{Journal of Artificial
  Intelligence Research}, 50, 105--140.
  \url{https://doi.org/10.1613/jair.4109}
\item
  Dubey, P., \& Shapley, L. S. (1979). Mathematical properties of the
  Banzhaf power index. \emph{Mathematics of Operations Research}, 4(2),
  99--131.
\item
  Felsenthal, D. S., \& Machover, M. (1998). \emph{The measurement of
  voting power: Theory and practice, problems and paradoxes}. Edward
  Elgar.
\item
  Kurz, S. (2015). Ready for the design of voting rules? \emph{Homo
  Oeconomicus}, 32(1), 117--129. (Preprint: arXiv:1405.0823).
\item
  Kurz, S. (2016). The inverse problem for power distributions in
  committees. \emph{Social Choice and Welfare}, 47(1), 65--88.
  \url{https://doi.org/10.1007/s00355-015-0946-8}
\item
  Kurz, S., \& Napel, S. (2014). Heuristic and exact solutions to the
  inverse power index problem for small voting bodies. \emph{Annals of
  Operations Research}, 215(1), 137--163.
\item
  Matsui, Y., \& Matsui, T. (2001). NP-completeness for calculating
  power indices of weighted majority games. \emph{Theoretical Computer
  Science}, 263(1--2), 305--310.
  \url{https://doi.org/10.1016/S0304-3975(00)00251-6}
\item
  Muroga, S. (1971). \emph{Threshold logic and its applications}. Wiley.
\item
  Nurmi, H. (1982). The problem of the right distribution of voting
  power. In M. J. Holler (Ed.), \emph{Power, Voting, and Voting Power}
  (pp.~203--212). Physica-Verlag (Springer).
\item
  Penrose, L. S. (1946). The elementary statistics of majority voting.
  \emph{Journal of the Royal Statistical Society}, 109(1), 53--57.
\item
  Taylor, A. D., \& Zwicker, W. S. (1999). \emph{Simple games:
  Desirability relations, trading, pseudoweightings}. Princeton
  University Press.
\end{itemize}

\end{document}
